\newcommand{\verfasser}{Phil Steinhorst}					% Verfasser des Dokuments
% % % % % % % % % % % % % % % % % % % % % % % % % % % % % % % % %
%		PhistGeneral.tex										%
%		General XeLaTeX-Settings For Math Documents				%
%		(based on 'mitschrift.tex' )							%
%																%
%		Author: Phil Steinhorst									%
% % % % % % % % % % % % % % % % % % % % % % % % % % % % % % % % %

\RequirePackage[thinlines]{easybmat} %-- muss aufgrund eines Bugs vor etex und tikz geladen werden
\documentclass[a4paper, twoside, headsepline, index=totoc,toc=listof, fontsize=9pt, cleardoublepage=empty, headinclude, DIV=13, BCOR=13mm, titlepage]{scrartcl}

\usepackage{scrtime} 	% KOMA, Uhrzeit ermoeglicht
\usepackage{scrpage2} 	% wie fancyhdr, nur optimiert auf KOMA-Skript, leicht andere Syntax
\usepackage{etoolbox}
\usepackage{ifxetex}

\usepackage[usenames, table, x11names]{xcolor}	% Farben. Vor TikZ laden!

% % % TikZ-Pakete % % %
\usepackage{tikz}
\usepackage{tikz-cd}
\usetikzlibrary{external}
\tikzset{>=latex}
\usetikzlibrary{shapes,arrows,intersections}
\usetikzlibrary{calc,3d}
\usetikzlibrary{decorations.pathreplacing,decorations.markings}
\usetikzlibrary{angles}
\tikzexternalize[prefix=tikz/, up to date check=diff]	% -shell-escape-Flag nötig!

% % % verwende LuaLaTeX, wegen dynamischer Speicherallokation
\tikzset{external/system call={xelatex \tikzexternalcheckshellescape 
    -halt-on-error -interaction=batchmode -jobname "\image" "\texsource"}}
    
% % % tikzexternalize für tikzcd deaktivieren, da inkompatibel
\AtBeginEnvironment{tikzcd}{\tikzexternaldisable}
\AtEndEnvironment{tikzcd}{\tikzexternalenable}

% % % um Inkompatibilitaeten von quotes und polyglossia bzw. babel zu vermeiden
\tikzset{
  every picture/.append style={
    execute at begin picture={\shorthandoff{"}},
    execute at end picture={\shorthandon{"}}
  }
}
\usetikzlibrary{quotes}
\usepackage{pgfplots}

% % % Mathe-Zeugs
\usepackage{mathtools} 						% beinhaltet amsmath
\mathtoolsset{showonlyrefs, centercolon}
\usepackage{wasysym}
\usepackage{amssymb} 						% zusätzliche Symbole
\usepackage{latexsym} 						% zusätzliche Symbole
\usepackage{stmaryrd} 						% für Blitz
\usepackage{nicefrac} 						% schräge Brüche
\usepackage{cancel} 						% Befehle zum Durchstreichen
\usepackage{extarrows}						% mehr Pfeile
\usepackage{mathdots}
\DeclareSymbolFont{bbold}{U}{bbold}{m}{n}
\DeclareSymbolFontAlphabet{\mathbbold}{bbold}
\newcommand{\mathds}[1]{\mathbb{#1}} 		% um Kompatibilitaet mit frueheren Benutzung von dsfont herzustellen
\def\mathul#1#2{\color{#1}\underline{{\color{black}#2}}\color{black}} %farbiges Untersteichen im Mathe-Modus

% % % XeTeX & Fonts
\usepackage{mathspec} 						% beinhaltet fontspec 
\usepackage{polyglossia} 					% babel-Ersatz
\setmainlanguage[spelling=new,babelshorthands=false]{german}
\newcommand\glqq{"}
\newcommand\grqq{"}
\defaultfontfeatures{Mapping=tex-text, WordSpace={1.4}} %
\setmainfont[Ligatures=Common, BoldFont={* Bold}, ItalicFont={* Light Italic}]{Source Sans Pro}
\setsansfont[Scale=MatchLowercase,Ligatures=Common, BoldFont={* Medium}]{Ubuntu}
\setallmonofonts[Scale=MatchLowercase, ItalicFont={*}]{Consolas} 
\usepackage{xltxtra}
\usepackage{fontawesome}

% % % Misc.
\usepackage[neverdecrease]{paralist}
\usepackage[german=quotes]{csquotes}
\usepackage{booktabs}
\usepackage{wrapfig}
\usepackage{float}
\usepackage[margin=10pt, font=small, labelfont=sf, format=plain, indention=1em]{caption}
\captionsetup[wrapfigure]{name=Abb. }
\usepackage{stackrel}
\usepackage{ifthen}
\usepackage{multicol}
%\flushbottom

% % % Unterstreichung
\usepackage[normalem]{ulem}
\setlength{\ULdepth}{1.8pt}

% % % Indexverarbeitung
\usepackage{makeidx}
\newcommand{\bet}[1]{\textbf{#1}}
\newcommand{\Index}[1]{\textbf{#1}\index{#1}}
\makeindex
\renewcommand{\indexpagestyle}{scrheadings}

% % % Marginnote & todonotes
\usepackage{marginnote}
\renewcommand*{\marginfont}{\color{gray} \footnotesize }
\usepackage[textsize=small]{todonotes}
\makeatletter
\renewcommand{\todo}[2][]{\tikzexternaldisable\@todo[#1]{#2}\tikzexternalenable}
\makeatother

% % % Konfiguration von Hyperref pdfstartview=FitH, 
\usepackage[hidelinks, pdfpagelabels,  bookmarksopen=true, bookmarksnumbered=true, linkcolor=black, urlcolor=RoyalBlue3, plainpages=false, hypertexnames=false, citecolor=black, hypertexnames=true, pdfauthor={\verfasser}, pdfborderstyle={/S/U}, linkbordercolor=RoyalBlue3, colorlinks=true, unicode, pdfencoding=auto]{hyperref}

\newcommand{\RM}[1]{\MakeUppercase{\romannumeral #1{}}} 	% Römische Zahlen
\usepackage{ellipsis}										% Punkte

% % % % % % % % % % % % % % % % % % % % % % % % % % % % % % % % %
%		PhistMath.tex											%
%		Weitere Mathe-Befehle									%
%																%
%		Author: Phil Steinhorst									%
% % % % % % % % % % % % % % % % % % % % % % % % % % % % % % % % %

% % % Buchstaben und Zahlen
\newcommand{\aff}{\mathbb{A}}
\newcommand{\CC}{\mathbb{C}}
\newcommand{\FF}{\mathbb{F}}
\newcommand{\HH}{\mathbb{H}}
\newcommand{\KK}{\mathbb{K}}
\newcommand{\NN}{\mathbb{N}}
\newcommand{\OO}{\mathbb{O}}
\newcommand{\QQ}{\mathbb{Q}}
\newcommand{\RR}{\mathbb{R}}
\newcommand{\ZZ}{\mathbb{Z}}
\newcommand{\bigO}{\mathcal{O}}					% Landau-O
\newcommand{\ind}{1\hspace{-0,8ex}1} 			% Indikatorfunktion (Doppeleins)

% % % Abk�rzungen
\newcommand{\ab}[1]{\overline{#1}}					% Abschluss
\newcommand{\assoz}{\mathrel{\hat=}}				% assoziiert
\newcommand{\bewrueck}{\glqq$\Leftarrow$\grqq:} 	% Beweis R�ckrichtung
\newcommand{\bewhin}{\glqq$\Rightarrow$\grqq:}		% Beweis Hinrichtung
\newcommand{\borel}{\mathfrak{B}}					% Borelsche Sigma-Algebra
\newcommand{\maxid}{\mathfrak{m}}					% maximales Ideal
\newcommand{\setone}{\{1\}}							% Einsmenge
\newcommand{\leb}{\lambda \hspace{-0,95ex}\lambda}	% Lebesgue-Ma� (Doppel-Lambda)
\newcommand{\Lp}{\mathcal{L}}						% L^p-R�ume
\newcommand{\NT}{\trianglelefteq}					% Normalteiler
\newcommand{\setnull}{\{0\}}						% Nullmenge
\newcommand{\weak}{\rightharpoonup}					% schwache Konvergenz
\newcommand{\weaks}{\overset{*}{\rightharpoonup}}	% schwache *-Konvergenz
\newcommand{\salg}{\mathfrak{A}}					% Sigma-Algebra (Skript-A)
\newcommand{\zyklot}[1]{#1^{(\infty)}}				% zyklotomische Erweiterung


% % % Operatoren
\DeclareMathOperator{\Alt}{Alt} 					% Alternierende n-Linearform
\DeclareMathOperator{\Aut}{Aut} 					% Automorphismen
\DeclareMathOperator{\Bil}{Bil} 					% Bilinearformen
\DeclareMathOperator{\bild}{Bild} 					% Bild
\DeclareMathOperator{\Char}{char} 
\DeclareMathOperator{\dom}{dom} 					% Domain
\DeclareMathOperator{\diam}{diam}					% Durchmesser
\DeclareMathOperator{\dist}{dist} 					% Distanz
\DeclareMathOperator{\eqs}{\mathrel{\widehat{=}}}	% entspricht
\DeclareMathOperator{\diver}{div} 					% Gradient
\DeclareMathOperator{\EPK}{EPK} 					% Einpunktkompaktifizierung
\DeclareMathOperator{\End}{End} 					% Endomorphismen
\DeclareMathOperator{\esssup}{esssup}				% essentielles Supremum
\DeclareMathOperator{\Gal}{Gal}	 					% Galoisgruppe
\DeclareMathOperator{\ggT}{ggT} 					% ggT
\DeclareMathOperator{\GL}{GL}						% allgemeine lineare Gruppe
\DeclareMathOperator{\grad}{grad} 					% Gradient
\DeclareMathOperator{\Grad}{Grad} 					% Grad
\DeclareMathOperator{\Hess}{Hess} 					% Hesse-Matrix
\DeclareMathOperator{\Hom}{Hom} 					% Homomorphismen
\DeclareMathOperator{\id}{id} 						% identische Abbildung
\DeclareMathOperator{\im}{im} 						% image
\DeclareMathOperator{\Jac}{Jac} 					% Jacobson-Radikal
\DeclareMathOperator{\Kern}{Kern}					% Kern
\DeclareMathOperator{\kgV}{kgV} 					% kgV
\DeclareMathOperator{\Koker}{Koker} 				% Kokern
\DeclareMathOperator{\Cov}{Cov} 					% Kovarianz
\DeclareMathOperator{\Mod}{Mod} 					% Moduln
\DeclareMathOperator{\modu}{mod} 					% Modulo
\DeclareMathOperator{\ord}{ord} 					% Ordnung
\DeclareMathOperator{\der}{\partial}				% Partielle Ableitung
\DeclareMathOperator{\pot}{\mathcal{P}}				% Potenzmenge
\DeclareMathOperator{\prlim}{\varprojlim\limits}	% projektiver Limes
\DeclareMathOperator{\Quot}{Quot}					% Quotientenring
\DeclareMathOperator{\Rang}{Rang} 					% Rang
\DeclareMathOperator{\rot}{rot} 					% Rotation
\DeclareMathOperator{\sgn}{sgn} 					% Signum
\DeclareMathOperator{\Spec}{Spec} 					% Spektrum
\DeclareMathOperator{\SL}{SL} 						% Spezielle lineare Gruppe
\DeclareMathOperator{\SO}{SO} 						% Spezielle orthogonale Gruppe
\DeclareMathOperator{\SU}{SU} 						% Spezielle unit�re Gruppe
\DeclareMathOperator{\Spur}{Spur} 					% Spur
\DeclareMathOperator{\supp}{supp} 					% Tr�ger
\DeclareMathOperator{\Sym}{Sym} 					% Symmetrische Gruppe
\DeclareMathOperator{\tr}{tr} 						% trace

% % % Klammerungen
\DeclarePairedDelimiter{\abs}{\lvert}{\rvert}		% Betrag
\DeclarePairedDelimiter{\ceil}{\lceil}{\rceil}		% aufrunden
\DeclarePairedDelimiter{\floor}{\lfloor}{\rfloor}	% aufrunden
\DeclarePairedDelimiter{\sprod}{\langle}{\rangle}	% spitze Klammern
\DeclarePairedDelimiter{\enbrace}{(}{)}				% runde Klammern
\DeclarePairedDelimiter{\benbrace}{\lbrack}{\rbrack}			% eckige Klammern
\DeclarePairedDelimiter{\penbrace}{\{}{\}}			% geschweifte Klammern

% % % Norm
\DeclarePairedDelimiter\doppelstrich{\Vert}{\Vert}
\newcommand{\norm}[2][\relax]{
\ifx#1\relax \ensuremath{\doppelstrich*{#2}}
\else \ensuremath{\doppelstrich*{#2}_{#1}}
\fi}							% selbst definierte Mathe-Befehle							% Grundeinstellungen (nicht ändern!)
\newcommand{\fach}{Partielle Differentialgleichungen I}	% Titel der Vorlesung								
\newcommand{\shortFach}{PDGL I}		% Titel kurz
\newcommand{\prof}{Prof.\,Dr.\,Benedikt Wirth}				% Dozent
\newcommand{\untertitel}{gelesen von \prof}					% Untertitel
\newcommand{\semester}{Sommersemester 2014}					% Semester
\newcommand{\homepage}{http://wwwmath.uni-muenster.de/num/Vorlesungen/PDEI_SS14/}	% Vorlesungshomepage
% % % % % % % % % % % % % % % % % % % % % % % % % % % % % % % % %
%		PhistScript.tex											%
%		XeLaTeX-Settings for Scripts							%
%																%
%		Author: Phil Steinhorst									%
% % % % % % % % % % % % % % % % % % % % % % % % % % % % % % % % %

\setlength\parindent{0pt}             % ohne Einrueckung

% % % Konfiguration von scrheadings
\setheadsepline{1pt}[\color{black}]
\pagestyle{scrheadings}
\clearscrheadfoot

% % % Kopf-/Fußzeilenlayout
\providecommand{\shortFach}{\fach}
\ohead{\small \leftmark}
\ofoot[{ \small \thepage}]{ \small \thepage} 
\automark{section}

% Metadaten
\title{\fach}
\author{\verfasser}
\subtitle{\untertitel}
\date{\semester}

% Inhaltsverzeichnis
\setcounter{tocdepth}{1}
\usepackage[tocindentauto]{tocstyle}
\usetocstyle{KOMAlike}	

% Nummerierung
\numberwithin{equation}{section}		% Section in Equation-Nummerierung einbeziehen

% ntheorem package
\newcounter{tcount}						% Zähler für Theoreme
\numberwithin{tcount}{section}			% Section in Theorem-Nummerierung einbeziehen
\usepackage[hyperref]{ntheorem}			% Lade ntheorem mit hyperref-Workaround
\theoremstyle{break}					% Zeilenumbruch nach Theoremkopf
\theorembodyfont{\normalfont}			% nicht kursiv
\theorempreskipamount0.5cm				% Abstand vor Theorem
\theorempostskipamount0.5cm				% Abstand nach Theorem
\newtheorem{thm}[tcount]{Theorem}
\newtheorem{bsp}[tcount]{Beispiel}
\newtheorem{defn}[tcount]{Definition}
\newtheorem{bem}[tcount]{Bemerkung}
\newtheorem{lemma}[tcount]{Lemma}

% Workaround für Tabulatoren im Theorem-Verzeichnis
\makeatletter
\def\thm@@thmline@name#1#2#3#4{%
        \@dottedtocline{-2}{0em}{2.3em}%
                   {\makebox[\widesttheorem][l]{#1 \protect\numberline{#2}}#3}%
                   {#4}}
\@ifpackageloaded{hyperref}{
\def\thm@@thmline@name#1#2#3#4#5{%
    \ifx\\#5\\%
        \@dottedtocline{-2}{0em}{2.3em}%
            {\makebox[\widesttheorem][l]{#1 \protect\numberline{#2}}#3}%
            {#4}
    \else
        \ifHy@linktocpage\relax\relax
            \@dottedtocline{-2}{0em}{2.3em}%
                {\makebox[\widesttheorem][l]{#1 \protect\numberline{#2}}#3}%
                {\hyper@linkstart{link}{#5}{#4}\hyper@linkend}%
        \else
            \@dottedtocline{-2}{0em}{2.3em}%
                {\hyper@linkstart{link}{#5}%
                  {\makebox[\widesttheorem][l]{#1 \protect\numberline{#2}}#3}\hyper@linkend}%
                    {#4}%
        \fi
    \fi}
}
\makeatother
\newlength\widesttheorem
\AtBeginDocument{
  \settowidth{\widesttheorem}{Proposition 10.10\quad}	% Referenzstring für die Breite der ersten Spalte
}

\theoremlisttype{optname}	% nur Theoreme auflisten, die einen Namen haben. 							% Seitenlayout für Skripte
\begin{document}
% % % % % % % % % % % % % % % % % % % % % % % % % % % % % % % % %
%		PhistTitle.tex											%
%		Layout für Skript-Titelseiten							%
%																%
%		Author: Phil Steinhorst									%
% % % % % % % % % % % % % % % % % % % % % % % % % % % % % % % % %
\begin{titlepage}
\pagestyle{empty}
\begin{center}
\begin{minipage}{0.4\textwidth}
\begin{flushleft}
\includegraphics[height=1.5cm,keepaspectratio]{../!config/wwulogo.pdf}
\end{flushleft}
\end{minipage}
\hfill
\begin{minipage}{0.4\textwidth}
\begin{flushright}
\vspace*{0.3cm}
\includegraphics[height=1.2cm,keepaspectratio]{../!config/fb10logo.pdf} \
\end{flushright}
\end{minipage}

% Title
\vspace*{6cm}
\textbf{\Huge{\fach}} \\
\vspace{0.2cm} 
\textbf{{\LARGE \untertitel}} \\
\vspace{0.6cm}
\LARGE{Mitschrift von \verfasser} \\
\vspace{0.6cm}
\LARGE{\semester} \\
\vspace*{2cm}
\footnotesize{\url{\homepage}}
\vfill

\vspace*{2cm}
\begin{flushright}
{\footnotesize Stand: \today}
\end{flushright}
\end{center}
\end{titlepage}							% Einstellungen für Titelseite
\section*{Vorwort}
\label{sec:preface}
	Der vorliegende Text ist eine inhaltliche Aufbereitung zur Vorlesung Elliptische Kurven und Kryptographie, gelesen von PD Dr. Karin Halupczok an der WWU Münster im Sommersemester 2015. Der Inhalt entspricht weitestgehend dem handschriftlichen Skript, welches auf der Vorlesungswebsite bereitsgestellt wird. Dieses Werk ist daher keine Eigenleistung des Autors und wird nicht von der Dozentin der Veranstaltung korrekturgelesen. Für die Korrektheit des Inhalts wird keinerlei Garantie übernommen. Bemerkungen, Korrekturen und Ergänzungen kann man folgenderweise loswerden:
	\begin{itemize}
		\item persönlich durch Überreichen von Notizen oder per E-Mail
		\item durch Abändern der entsprechenden TeX-Dateien und Versand per E-Mail an mich
		\item direktes Mitarbeiten via GitHub. Dieses Skript befindet sich im \texttt{latex-wwu}-Repository von Jannes Bantje:
		\begin{center}
			\url{https://github.com/JaMeZ-B/latex-wwu}
		\end{center}
	\end{itemize}

\subsection*{Literatur}
\label{sub:lit}
\begin{itemize}
	\item Blake, Seroussi, Smart: Elliptic curves in cryptography
	\item Menezes, van Oorschot, Vanstone: Handbook of applied cryptography
	\item Silverman: The arithmetic of elliptic curves
	\item Silverman: A friendly introduction to number theory, chap. 40-45
	\item Washington: Elliptic curves, number theory and cryptography
	\item Werner: Elliptische Kurven in der Kryptographie
\end{itemize}

\subsection*{Kommentar der Dozentin}
In der Vorlesung beschäftigen wir uns mit den arithmetischen und geometrischen Eigenschaften elliptischer Kurven sowie deren Anwendungen in der Kryptographie. Dabei werden wir auch einen Vergleich mit Anwendungen der elementaren Zahlentheorie in der Kryptographie ziehen. Wir verfolgen eine elementare Herangehensweise, d.h. Kenntnisse der algebraischen Geometrie und der Funktionen- oder Zahlentheorie werden nicht benötigt. Es genügen die Vorkenntnisse aus den Grundvorlesungen.

\subsection*{Vorlesungswebsite}
\label{sub:link}
Das handgeschriebene Skript sowie weiteres Material findet man unter folgendem Link:
\begin{center}
	\url{\homepage}
\end{center}

\subsection*{Titelbild}
\label{sub:titlepic}
Das fehlt noch. Über Ideen und Anregungen freue ich mich sehr!

\vfill
\begin{flushright}
	Phil Steinhorst \\
	p.st@wwu.de
\end{flushright}
\newpage
\tableofcontents					% Inhaltsverzeichnis ausgeben
\newpage
% % % % % % % % % % % % % % % % % % % % % % % % % % % % %
\chapter{CAT(0)-Räume} % (fold)
\label{cha:1}

\begin{definition}[metrischer Raum]
	Sei $X$ eine Menge. Eine Abbildung $d \colon X \times X \rightarrow \RR_{\geq 0}$ heißt \Index{Metrik}, wenn für alle $x,y,z \in X$ gilt:
	\begin{enumerate}[(i)]
		\item $d(x,y) = 0 \Leftrightarrow x=y$
		\item $d(x,y) = d(y,x)$
		\item $d(x,z) \leq d(x,y) + d(y,z)$
	\end{enumerate}
	Das Paar $(X,d)$ heißt dann \Index{metrischer Raum}.
\end{definition}

\cleardoubleoddemptypage				% Einleitung
\section{Der euklidische Algorithmus}
\label{sec:para2}
	Sei $R$ kommutativer Ring mit $1 \neq 0$. Für beliebiges $a \in R$ betrachte man die Menge der Vielfachen von $a \in R$, also
	\[ Ra := {xa : x \in R} = {b \in R : a|b} \]
	Die Teilmenge $I = Ra$ hat folgende Eigenschaften:
	\begin{enumerate}[(i)]
		\item $0 \in I$
		\item $b_1,b_2 \in I \Rightarrow b_1+b_2 \in I$
		\item $c \in R, b \in I \Rightarrow cb \in I$
	\end{enumerate}
	
\begin{defn}[Ideal, Hauptideal] \label{def_2.1}
	Eine Teilmenge $I$ von $R$ heißt ein \Index{Ideal} in $R$, falls die Eigenschaften (i), (ii), (iii) erfüllt sind. $I$ heißt \Index{Hauptideal}, wenn es ein $a \in R$ gibt mit $I = Ra$. Wir verwenden die Bezeichnung
	\[ (a) := Ra \]
	und nennen $(a)$ das von $a \in R$ erzeugte Hauptideal.
\end{defn}

\minisec{Bemerkung}
	\begin{enumerate}[(1)]
		\item $(b) \subseteq (a) \Leftrightarrow a | b$
		\item $a \assoz b \Leftrightarrow (a) = (b)$
		\item $c$ ist gemeinsames Vielfaches von $a_1,\dots,a_n \Leftrightarrow (c) \subseteq (a_1) \cap \dots \cap (a_n)$
		\item $m$ ist ein kgV von $a_1,\dots,a_n \Leftrightarrow (a_1) \cap \dots \cap (a_n) = (m)$
		\item $d$ ist ein gemeinsamer Teiler von $a_1,\dots,a_n \Leftrightarrow (a_i) \subseteq (d)$ für $1 \leq i \leq n$
		\item $d$ ist ein gemeinsamer Teiler von $a_1,\dots,a_n \Leftrightarrow Ra_1 + Ra_2 + \dots + Ra_n \subseteq (d)$
		\item $d$ ist ein ggT von $a_1,\dots,a_n \Leftrightarrow (d)$ ist das kleinste Hauptideal mit $Ra_1 + \dots Ra_n \subseteq (d)$.
	\end{enumerate}
	Ein ggT lässt sich also idealtheoretisch nicht so einfach charakterisieren wie oben ein kgV durch (4). Am schönsten wäre es, wenn $Ra_1 + \dots + Ra_n$ ein Hauptideal wäre, dann würde (7) übergehen in:
	\[ d \text{ ist ein ggT von } a_1,\dots,a_n \Leftrightarrow Ra_1+Ra_2+ \dots + Ra_n = (d) \]
	
\begin{defn}[Hauptidealring] \label{def_2.2}
	Ein Integritätsring $R$ heißt ein \Index{Hauptidealring}, wenn jedes Ideal $I$ von $R$ ein Hauptideal ist.
\end{defn}

\minisec{Bezeichnung}
	Für Elemente $a_1,\dots,a_n$ in einem beliebigen kommutativen Ring $R$ mit $1 \neq 0$ setze
	\[ (a_1,\dots,a_n) := Ra_1 + \dots + Ra_n \]
	Man nennt $(a_1,\dots,a_n)$ das von $a_1,\dots,a_n$ erzeugte Ideal in $R$.
\newpage				% Cauchy-Kovalevskaya
\section{Partielle Differentialgleichungen erster Ordnung}
\label{sec:para3}

% % % % % % 22. Apr.
\minisec{Motivation: Linearer Fall}
	Betrachte \marginnote{22. Apr}
	\[ a(x) \cdot \nabla u(x) + a_0(x) = 0 \] mit Randdaten auf glatter $(n-1)$-dimensionaler Hyperfläche $\Gamma \subset \RR^n$, $u(x) = g(x)$ auf $\Gamma$. \\
	Die Richtungsableitung von $u$ in Richtung $a$ ist $a(x) \cdot \nabla u(x) = -a_0 (x)$, d.h. wir wissen, wie sich $u$ entlang charakteristischer Kurven $s \mapsto x(s)$ ändert mit
	\begin{equation}
	\begin{aligned}
		\dot{x}(s) = a(x(s)) \label{eq_2204_1}
	\end{aligned}
	\end{equation} \\
	Sei $\Gamma$ parametrisiert durch $x_0 \colon \RR^{n-1} \rightarrow \Gamma$. Zu $\widehat{x} \in \RR^n$ können wir $s \in \RR, y \in \RR^{n-1}$ finden, sodass $\widehat{x} = x(s)$ für eine Lösung $x(s)$ von \eqref{eq_2204_1} mit $x(0) = x_0(y)$, d.h. wir reparametrisieren $\RR^n$ durch eine Funktion $\widehat{x}\colon \RR^{n-1} \times \RR \rightarrow \RR, (y,s) \mapsto \widehat{x}(y,s)$. Dann gilt:
	\[ u(\widehat{x}(y,s)) = \underbrace{u(\widehat{x}(y,0))}_{g(\widehat{x}(y,0))} + \int\limits_{0}^{s} \underbrace{\frac{\der u(\widehat{x}(y,t))}{\der \tau}}_{=-a_0(\widehat{x}(y,\tau))} d\tau; \quad u(x) = (u \circ \widehat{x}) \circ \widehat{x}^{-1} \]
	
\begin{bem} \label{bem_2}
	Damit $u$ auf $\RR^n$ definiert ist, muss $\widehat{x}^{-1}$ definiert sein.\marginnote{[2]} Für lokale Invertierbarkeit ist hinreichend:
	\[ \det(D \widehat{x}) \neq 0 \]
	Auf $\Gamma$ ist $\det(D \widehat{x}(y,0) = \det( \underbrace{D x_0(y)}_{\in R^{n\times m -1}}\vert a(x_0(y)))$, d.h. wir brauchen $a(x_0(y)) \notin \operatorname{span}(Dx_0(y))$. Die Spalten von $Dx_0(y)$ spannen den Tangentialraum von $\Gamma$ an $x_0(y)$ auf. Damit dies ungleich 0 ist, muss $a \notin \operatorname{span}(Dx_0)$, d.h. $a \cdot \nu \neq 0$ mit $\nu$ als Normale auf $\Gamma$, d.h. $\Gamma$ ist nichtcharakteristisch.
\end{bem}
	
\begin{bsp} \label{bsp_3}
	Betrachte \marginnote{[3]}
	\[ \begin{cases}
		u_{x_1} + u_{x_2} = 1 \text{ auf } \RR^2 \\
		u = 0 \text{ auf } \Gamma = \{x_1 + x_2 = 0 \}
	\end{cases} \]
	Charakteristische Kurven: $\dot{x} = \begin{pmatrix} 1 \\ 1	\end{pmatrix} \Rightarrow x(s) = \begin{pmatrix} c_1 + s \\ c_2 + s \end{pmatrix}$. \\
	Parametrisiere $\Gamma$:
	\[ x_0(y) := \begin{pmatrix} y \\ -y \end{pmatrix} = x_0(y) \]
	$\Rightarrow \widehat{x}(y,s) = x(s)$ für die Kurve $x(s)$ mit $x(0) = x_0(y) \Rightarrow \widehat{x}(y,s) = \begin{pmatrix} y+s \\ -y+s \end{pmatrix}$ \\
	$\Rightarrow \widehat{x}^{-1}(x) = (\frac{x_1-x_2}{2},\frac{x_1+x_2}{2})$ \\
	$\Rightarrow u(\widehat{x}(y,s)) = 0 + \int\limits_{0}^{s} 1 d\tau = s$ \\
	$\Rightarrow u(x) = \underbrace{(u \circ \widehat{x})}_s \circ \widehat{x}^{-1}(x) = \frac{x_1+x_2}{2}$
\end{bsp}
	
\minisec{Nichtlinearer Fall}
	Betrachte
	\begin{equation}
		0 = F(\nabla u(x),u(x),x), \quad u=g \text{ auf } \Gamma, \quad F,g,\Gamma \ C^2\text{-glatt} \label{eq_2204_2}
	\end{equation}
	Idee: Information fließt entlang charakteristischer Kurven \\
	Ziel: Finde charakteristische Kurven und wandele die PDGL um in ein System von gewöhnlichen DGL entlang der Kurven. \\
	Notation: \begin{itemize}
	\item Für eine Kurve $x(s)$ schreibe $z(s) := u(x(s))$ und $p(s) = \nabla u(x(s))$, d.h.
	\[ 0 = F(p(s), z(s), x(s)) \]
	\item $\dot{} = \frac{\der}{\der s}$
	\item $F_p,F_z,F_x$ sind partielle Ableitungen von $F$ nach erstem, zweitem und drittem Argument.
	\end{itemize}
	Suche GDGL für $x,p,z$.
	
\begin{thm}[Charpits Gleichungen]
\label{thm_charpits} \label{thm_4}
	Sei $u \in C^2(\Omega)$ eine Lösung von \eqref{eq_2204_2}.\marginnote{[4]} Entlang von Kurven $x(s)$ in $\Omega$ mit
	\begin{equation}
	\begin{aligned}
		\dot{x} = F_p \label{eq_2204_3}
	\end{aligned}
	\end{equation}
	gilt
	\begin{align}
		&\dot{z} = p \cdot F_p \label{eq_2204_4} \\
		&\dot{p} = -pF_z - F_x \label{eq_2204_5}
	\end{align}
	Diese GDGLs heißen \Index{Charpits} oder \Index{charakteristische Gleichungen}, Kurven $x(s)$ oder $(x(s),z(s),p(s))$ heißen \Index{charakteristische Kurven}.
\end{thm}
	
\minisec{Beweis}
	Betrachte beliebige Kurve $x(s)$. Dann gilt:
	\[ \begin{array}{c}
		\dot{z}(s) = \nabla u(x(s)) \cdot \dot{x}(s) = p(s) \cdot \dot{x}(s) \\
		\dot{p}(s) = D^2u(x(s)) \dot{x}(s)
	\end{array} \]
	Um $D^2u$ loszuwerden, differenziere \eqref{eq_2204_2} nach $x$:
	\[ 0 = \frac{d}{dx} F(\nabla u(x),u(x),x) = D^2 u F_p + F_z \overbrace{\nabla u}^p + F_x \]
	Wähle die Kurve $x(s)$, sodass $\dot{x}(s) = F_p(p(s),z(s),x(s))$, so erhalten wir
	\[ \dot{p}(s) = -F_z p - F_x \qquad \dot{z} = p \cdot F_p \]. \qed

\mbox{} \\
\eqref{eq_2204_3}, \eqref{eq_2204_4} und \eqref{eq_2204_5} können mit Anfangswerten $x(0),z(0),p(0)$ gelöst werden. Sei $\Gamma$ parametrisiert durch $x_0(y), y \in \RR^{n-1}$. Werte für $x(0)$ und $z(0)$ sind auf $\Gamma$ gegeben durch $x(0) = x_0(y) \in \Gamma$ und $z(0) = g(x_0(y))$. $p(0)$ erhalten wir durch Lösen von
\begin{equation}
	\left. \begin{cases}
	0 = F(p(0),z(0),x(0)) & \text{PDGL} \\
	0 = \nabla x_0 (y) p(0) - \nabla_y g(x_0(y)) & \text{Abl. von } u(x_0(y)) = g(x_0(y)) \text{ nach } y
	\end{cases} \right\}  \label{eq_2204_7} \end{equation}
wobei
	\[ \nabla x_0 (y) = \begin{pmatrix}
	\der_{y_1} {x_0}_1 & \cdots & \der_{y_1} {x_0}_n \\ 
	\vdots & \ddots & \vdots \\ 
	\der_{y_{n-1}} {x_0}_1 & \cdots & \der_{y_{n-1}} {x_0}_n
	\end{pmatrix} \]
Die Lösung von \eqref{eq_2204_3}, \eqref{eq_2204_4} und \eqref{eq_2204_5} für diese Anfangswerte liefert eine Parametrisierung $\widehat{x}(y,s)$ des $\RR^n$ mit $\widehat{x}(y,s) = x(s)$ für die Kurve $x(s)$ mit $x(0) = x_0(y)$ sowie Funktionen $z(y,s),p(y,s)$.

\begin{defn}[nichtcharakteristische Hyperfläche] \label{def_nichtchar} \label{def_5}
	Eine glatte $(n-1)$-dimensionale Hyperfläche\marginnote{[5]} $\Gamma \subset \RR^n$ heißt \Index{nichtcharakteristisch} zu \eqref{eq_2204_2}, wenn für die Einheitsnormale $\nu$ an $\Gamma$ gilt:
	\[ F_p(p,z,x) \cdot \nu(x) \neq 0 \text{ für alle } x \in \Gamma, z \in \RR, p \in \RR^n \]
\end{defn}
	
\begin{lemma} \label{lemma_6}
	Wenn gilt: \marginnote{[6]} \begin{itemize}
	\item $\Gamma$ ist nichtcharakteristisch
	\item Für ein gegebenes $y_0 \in \RR^{n-1}$ existiert eine Lösung von \eqref{eq_2204_7}
	\end{itemize}
	dann: \begin{itemize}
		\item kann \eqref{eq_2204_7} auch in einer Umgebung $U$ von $y_0$ nach $p(0)$ gelöst werden
		\item existieren ein offenes Intervall $I = (-\delta,\delta)$ und eine Umgebung $V$ von $x_0(y_0)$, sodass für alle $x \in V$ ein $y \in U$ und $s \in I$ existiert mit $x = \widehat{x}(y,s)$. $\widehat{x}^{-1}$ ist $C^2$.
	\end{itemize}
\end{lemma}
	
\minisec{Beweis}
	\begin{itemize}
		\item Existenz von $p(0)$ folgt direkt aus dem Satz über implizite Funktionen. (Beachte:
		\[ \frac{\der}{\der p(0)} \begin{pmatrix}
		F(p(0),z(0),x(0)) \\ \nabla x_0 (y) p(0) - \nabla_y g(x_0(y)) \end{pmatrix} = \begin{pmatrix}
		F_p \\ \nabla x_0 (y) \end{pmatrix} \]
		hat vollen Rang.)
		\item Existenz und Glattheit von $\widehat{x}^{-1}$ folgt aus dem Satz über inverse Funktionen, da $\det (y_0,0) = \det \begin{pmatrix}	\nabla x_0 \\ F_p \end{pmatrix} \neq 0$ \qed
	\end{itemize}
	
\begin{bem} \label{bem_8}
	Auch hier bedeutet nichtcharakteristisch,\marginnote{[7]} dass auf $\Gamma$ nach den partiellen Ableitungen von $u$ gelöst werden kann.
\end{bem}
	
\begin{thm}[Lokale Existenz einer glatten Lösung]
\label{thm_lok_existenz} \label{thm_8}
	Seien \marginnote{[8]} \begin{itemize}
		\item $\Gamma, g, F$ glatt, also $C^2$
		\item $\Gamma$ nichtcharakteristisch und
		\item es gebe ein $y_0 \in \RR^{n-1}$ und ein $p(0) \in \RR^n$, das \eqref{eq_2204_7} löst.
	\end{itemize}
	Dann existiert auf der Umgebung $V$ von $x_0(y_0)$ aus Lemma \ref{lemma_6} eine $C^2$-Lösung $u$ von \eqref{eq_2204_2} und \eqref{eq_2204_3}. Es ist $u(x) = z(y,s)$ und $\nabla u(x) = p(y,s)$ für $(y,s) = \widehat{x}^{-1} (x)$.
\end{thm}

% % % % % % % % % % % % % 25.04.
\minisec{Beweis}
	\begin{enumerate}[1.]
		\item Nach Lemma \eqref{lemma_6} \marginnote{25. Apr} existiert $\widehat{x}^{-1}$ auf $V$ sowie $z(y,0) = g(x_0(y)),p(y,0)$ mit $F(p(y,0),z(y,0),\widehat{x}(y,0)) = 0$ und $\nabla x_0 (y) p(0) - \nabla_y g(x_0(y)) = 0$. Löse Charpits Gleichungen mit diesen Anfangswerten $\Rightarrow z(y,s),p(y,s)$.
		\item Zeige: $f(y,s) := F(p(y,s),z(y,s),\widehat{x}(y,s)) = 0$. \\
			$f(y,0) = 0, \dot{f}(y,s) = F_p \cdot \dot{p} + F_z \cdot \dot{z} + F_x \dot{\widehat{x}} = F_p(-pF_z - F_x) + F_zpF_p + F_x \cdot F_p = 0$
		\item Zeige $p(\widehat{x}^{-1}(x)) = \nabla u(x)$ für $u(x) = z(\widehat{x}^{-1}(x))$. \marginnote{$z \equiv u \circ \widehat{x}^{-1}$ \\ $z_s = \nabla u \cdot \widehat{x}_s = p \cdot \widehat{x}_s$ \\ $z_y = \nabla u \widehat{x}_y = p \cdot \widehat{x}_y$} \begin{itemize}
				\item $\cdot{z}(y,s) = p(y,s) \cdot \dot{\widehat{x}}(y,s)$
				\item $z_y(y,s) = p^T(y,s) \widehat{x}_y(y,s)$, denn $r(s) := z_y(y,s) - p^T(y,s) \widehat{x}_y (y,s)$ erfüllt
				\begin{equation}
				\begin{aligned}
					\dot{r}(s) &= \underbrace{\dot{z_y}(y,s)}_{p^T_y \dot{\widehat{x}} + p^T \dot{\widehat{x}}_y} - \dot{p}^T(y,s) \widehat{x}_y(y,s) - p^T(y,s) \dot{\widehat{x}}_y(y,s) = p_y^T F_p - (-pF_z - F_x)^T \widehat{x}_y \\ \notag
					&= \frac{\der}{\der y} \underbrace{F(p(y,s),z(y,s),\widehat{x}(y,s))}_{=0} + \underbrace{p^T F_z \widehat{x}_y - F_z z_y}_{-F_z r}
				\end{aligned}
				\end{equation} 
				mit Anfangswert $r(0) = \underbrace{z_y(y,0)}_{\nabla_y g(x_0(y))} - p^T(y,0) \underbrace{\widehat{x}_y(y,0)}_{\nabla_y x_0(y)} = 0$. $\Rightarrow r \equiv 0$.
				\item $u_x = z_y y_x + z_s s_x = p^T \widehat{x}_y y_x + p^T \widehat{x}_s s_x = p^T \underbrace{(\widehat{x}_y y_x + \widehat{x}_s s_y)}_{\widehat{x}_x} = p^T, p = \nabla u$. \qed
			\end{itemize}
	\end{enumerate}

\begin{bsp} \label{bsp_9}
	Löse das Cauchy-Problem \marginnote{[9]}
	\[ \begin{cases}
		u_t + uu_x = 1 & \text{ auf } (0,\infty) \times \RR \\
		u = x & \text{ auf } \Gamma = \{t = 0\} \end{cases} \]
	Charakteristische Kurven sind gegeben durch $\frac{\der}{\der s} \begin{pmatrix} (t,x) \\ z \\ 	p \end{pmatrix} = \begin{pmatrix} (1,z) \\ P-1+zp_2 \\ -p_2 p \end{pmatrix}$, also
	\[ \begin{pmatrix} (t,x) \\ z \\ 	p \end{pmatrix} = \begin{pmatrix} (t(0)+s,x(0) + z(0)s + \frac{p_1(0)+z(0)p_2(0)}{2} s^2 \\ p_1(0)s + z(0)(1+p_2(0)s) \\ \frac{p}{1+p_2(0)s} \end{pmatrix}.\]
	Parametrisiere $\Gamma$ durch $x_0(y) = (0,y)$. Das liefert Anfangswerte $(t(0),x(0))=x_0(y)=(0,y)$ und $z(0) = y$. Lösen von \eqref{eq_2204_7} liefert $p(0) = \begin{pmatrix} 1-y \\ 1 \end{pmatrix}$. Insgesamt ergibt sich:
	\[ \Rightarrow \begin{pmatrix} (t,x) \\ z \\ p \end{pmatrix} = \begin{pmatrix} s,y+ys+\frac{s^2}{2} \\ y+s \\ \frac{1}{1+s} \begin{pmatrix} 1-y \\ 1 \end{pmatrix} \end{pmatrix} \Rightarrow (s,y) = \left( t,\frac{x-\frac{t^2}{2}}{1+t}\right) \Rightarrow u(t,x) = z(y,s) = y+s = t + \frac{x-\frac{t^2}{2}}{1+t} \]
\end{bsp}
	
\begin{bem} \label{bem_10}
	Für quasilineare PDGL der Form \marginnote{[10]}
	\[ b(u(x),x) \cdot \nabla u(x) + c(u(x),x) = 0 \]
	können die charakteristischen Gleichungen \eqref{eq_2204_3}, \eqref{eq_2204_4} und \eqref{eq_2204_5} vereinfacht werden: $b(z(s),x(s)) \cdot p(s) + c(z(s),x(s)) = 0$ impliziert:
	\begin{align}
		\dot{x}(s) = b(z(s),x(s))	\label{eq_2504_8} \\
		\dot{z}(s) = -c(z(s),x(s))	\label{eq_2504_9} \\
		\dot{p}(s) = -(b_z \cdot p + c_z)\cdot p - b_x p -c_x	\label{eq_2504_10}
	\end{align}
	Die ersten beiden Gleichungen legen bereits $x$ und $z$ fest, sodass \eqref{eq_2504_10} ignoriert werden kann.
\end{bem}
\newpage				% PDGL 1. Ordnung
\section{Sichere Kryptographie mit elliptischen Kurven}
\label{sec:para4}
\nextlecture
\subsection{Bekannte Angriffe auf das DL-Problem: Überblick}
\label{sub:4.1}
\begin{bem}
	Die Sicherheit des ElGamal- und ECDSA-Verfahrens beruht hier auf der Schwierigkeit des DL-Problems auf elliptischen Kurven. \marginnote{[19]}
	Allerdings gibt es bestimmte Arten elliptischer Kurven, bei denen das DL-Problem algorithmisch schnell lösbar ist, sodass sich diese Kurven als kryptographisch schwach bzw. ungeeignet erweisen.
	Auch die Wahl eines Punktes $P$ mit großer Ordnung ist wichtig.
\end{bem}

\subsubsection{BSGS und Silver-Pohlig-Hellman}
\label{subsub:4.1.1}
	Diese beiden Methoden eignen sich zur Lösung des DL-PRoblems in einer beliebigen abelschen Gruppe $G$.
	
\begin{defn}[DL-Problem in $G$]
	Gegeben sei $P \in G$ mit $\ord(P) = n \in \NN$, sowie $Q \in \sprod{P}$. Gesucht ist $k \in \{0,\dots,n-1\}$ mit $kP = Q$.
\end{defn}

\begin{defn}[BSGS -- Baby Steps Giant Steps]
	Dieses Verfahren kommt in Frage, wenn $P$ kleine Ordnung $n$ hat.
	Dann kann das DL-Problem wie folgt gelöst werden; der Algorithmus hat einen Zeit- und Platzbedarf der Größenordnung $\oh(\sqrt{n})$: \index{BSGS (Baby Steps Giant Steps)}
\end{defn}

\begin{bem}
	Vorüberlegung: Sei $m = \ceil*{\sqrt{n}} = \min \{ l \in \NN : l \geq \sqrt{n}\}$, schreibe $k = qm + r, r \in \{0,1,\dots,m-1\}$ (Division mit Rest). \\
	Ziel: Bestimme $q, r$. \\
	Da $Q = kP = qmP + rP$, folgt $\underbrace{Q - rP}_{\text{"Baby step"}} = \underbrace{qmP}_{\text{"Giant step"}}$.
\end{bem}

\begin{bem}
	Idee: Berechne alle möglichen Werte von "Baby step" und nach und nach die möglichen Werte von "Giant step".
	Trifft man auf eine Übereinstimmung, sind $r$ und $m$ gefunden.
\end{bem}

\begin{anw}[Schritt 1]
	Berechne die Liste der "Baby steps" $B = \{(Q-rP,r) : 0 \leq r \leq m \}$.
\end{anw}

\begin{anw}[Schritt 2]
	\begin{itemize}
		\item Ist für eines der $r$ die Gleichung $Q - rP = \oh$ erfüllt, ist $k = r$.
		\item Sonst teste für den ersten "Giant step" $R = mp$, ob $R$ in $B$ schon vorkommt.
		Falls ja: $k = m + r$.
		\item Teste so alle "Giant steps" $2R, 3R, 4R, \dots, (m-1)R$, ob diese in $B$ vorkommt.
		Wenn ja, gibt die zweite Komponente $r$ mit $k = qm+r$.
	\end{itemize}
\end{anw}

\begin{defn}[Silver-Pohlig-Hellman-Verfahren]
	Dieses Verfahren löst das DL-Problem in einer abelschen Gruppe $G$, wenn die Ordnung $n = \ord(P)$ aus nur kleinen Primfaktoren $p_i$ zusammengesetzt ist, d.h. \Index{glatt} ist. \index{Silver-Pohlig-Hellman}
\end{defn}

\begin{defn}[$B$-glatt]
	Sei $B \in \RR_{>0}$.
	Dann heißt $m \in \NN$ \bet{$B$-glatt}, falls für alle $p \mid n$ gilt: $p \leq B$.
	\index{B-glatt@$B$-glatt}
\end{defn}

\begin{bem}
	Die Bestimmung von $k$ in $\sprod{P}$ wird auf Untergruppen von $\sprod{P}$ der Ordnungen $p_i \mid n$ zurückgeführt.
	Sei $\ord(P) = n = \prod_{i=1}^{t} p_i^{\lambda_i}$ mit $p_1, \dots, p_t$ paarweise verschieden und prim, $\lambda_i \in \NN$.
	Der Algorithmus hat dann eine Laufzeit von $\oh \enbrace*{\sum_{i=1}^{t} (\lambda_i(\log n + \sqrt{p_i}))}$.
\end{bem}

\begin{bem}
	Vorüberlegung:
	\begin{itemize}
		\item Zur Bestimmung von $k$ mit $kP = Q \in \sprod{P}$ berechnen wir alle Restklassen $k \modu p_1^{\lambda_1}, k \modu p_2^{\lambda_2}, \dots, k \modu p_t^{\lambda_t}$.
		Denn laut chinesischem Restsatz ist $\ZZ/n\ZZ \simeq \ZZ/p_1^{\lambda_1} \ZZ \times \dots \times \ZZ / p_t^{\lambda_t} \ZZ$, sodass damit dann auch die Restklasse von $k \modu n$ bestimmt werden kann.
		\item Betrachte daher jedes $p = p_i, \lambda = \lambda_i$ mit $1 \leq  i \leq t$. \\
		Gesucht: $z \in \{0, \dots, p^\lambda - 1\}$ mit $z \kon k \modu p^\lambda$. \\
		Schreibe $z = z_0 + z_1 p + \dots + z_{\lambda-1} p^{\lambda-1}$, die $z_i \in \{0, \dots, p-1\}$, in der $p$-adischen Entwicklung; bestimme die $z_0, \dots, z_{\lambda-1}$.
	\end{itemize}
\end{bem}

\begin{anw}[Schritt 1]
	Sei $R := \frac{n}{p} \cdot P$, dann ist $\frac{n}{p}Q = \frac{n}{p}kP = kR$ und $pR = \oh$.
	Also ist $kR = zR = z_0 R$, d.h. $z_0 R = \frac{n}{p} \cdot Q$.
	Somit muss man in der Untergruppe $\sprod{R}$ der (kleinen) Ordnung $p$ ein DL-Problem lösen, um $z_0$ zu bestimmen -- etwa mit BSGS.
\end{anw}

\begin{anw}[Schritt 2]
	Seien $z_0, \dots, z_{j-1}$ schon (rekursiv) bestimmt, wo $j \leq \lambda-1$ ist.
	Berechne dann $Q_j := \frac{n}{p^{j+1}} (Q - (z_0 + z_1p + \dots + z_{j-1} p^{j-1}) P)$. Da $nP = \oh$, ist $\frac{n}{p^{j+1}} \cdot p^\lambda P = \oh$; da $z \kon k \modu p^\lambda$, ist $k = z + sp^\lambda, s \in \ZZ$, also $\frac{n}{p^{j+1}} Q = \frac{n}{p^{j+1}} k P = \frac{n}{p^{j+1}} zP + \underbrace{\frac{n}{p^{j+1}} \cdot sp^\lambda P}_{=\oh} = \frac{n}{p^{j+1}} zP$ und somit $Q_j = \frac{n}{p^{j+1}} (z_j p^j + \dots + z_{\lambda - 1} p^{\lambda-1}) P = \frac{n}{P} z_j P = z_j R$. \\
	Zur Berechnung von $z_j$ ist wieder ein DL-Problem in der Untergruppe $\sprod{R}$ der Ordnung $p$ zu lösen -- etwa mit BSGS.
\end{anw}

\begin{bem}
	Ist die Gruppenordnung glatt, ist der Algorithmus also sehr schnell.
\end{bem}

\subsubsection{Pollard-$\rho$ und Pollard-$\lambda$}
\label{subsub:4.1.2}
\begin{bem}
	Der \bet{Pollard-$\rho$-Algorithmus} ist von der Laufzeit her vergleichbar mit BSGS, ist aber speicherplatztechnisch günstiger und lässt sich gut parallelisieren.
	Mit $m$ Prozessoren wird der Algorithmus so um den Faktor $m$ schneller. \index{Pollard-$\rho$}
\end{bem}

\begin{bem}
	Der \bet{Pollard-$\lambda$-Algorithmus} ist ähnlich, im Allgemeinen aber eher langsamer als Pollard-$\rho$.
	Er liefert gute Ergebnisse, wenn der diskrete Logarithmus in einem hinreichend kleinen Intervall liegt.
	Auch Pollard-$\lambda$ ist gut parallelisierbar.
	(Die genauen Verfahren können in der Fachliteratur nachgeschlagen werden.) \index{Pollard-$\lambda$}
\end{bem}

\subsubsection{MOV und SSSA}
\label{subsub:4.1.3}
\begin{bem}
	Beim \Index{MOV-Verfahren} (Autoren: Menezes, Okamoto, Vanstone) wird das DL-Problem für eine elliptische Kurve $E(\FF_{p^r})$ auf das in der Gruppe $(\FF_{p^{rl}}^*,\cdot)$ für ein $l \geq 1$ zurückgeführt.
	Es ist also speziell nur für elliptische Kurven konstruiert, nicht für allgemeine abelsche Gruppen.
	Zeigt sich hier, dass $l \geq 1$ so wählbar ist, dass das DL-Problem in $(\FF_{p^{rl}}^*,\cdot)$ leicht, d.h. schnell, zu lösen ist, ist die elliptische Kurve kryptographisch ungeeignet, etwa wenn $n = \ord(P)$ Teiler von $p^{rl} - 1$ ist.
\end{bem}

\begin{bem}
	Generell lässt sich das DL-Problem in $(\FF_{p^{rl}}^*,\cdot)$ in subexponentieller Zeit schnell lösen (mit so genannten Indexkalkül-Methoden), sodass Kurven, für die das DL-Problem auf ein schnelles in einem $(\FF_{p^{rl}}^*,\cdot)$ zurückgeführt werden kann, als kryptographisch schwach bzw. ungeeignet angesehen werden.
	Das ist etwa bei supersingulären elliptischen Kurven der Fall, bei denen die Gruppenstruktur recht gut bekannt ist.
\end{bem}

\begin{defn}[supersingulär]
	Eine elliptische Kurve $E(\FF_{p^r})$ heißt \Index{supersingulär}, falls $ p =\Char(\FF_{p^r})$ die Spur des Frobenius teilt, d.h. $p \mid p^r +1 - \#E(\FF_{p^r})$.
\end{defn}

\begin{bem}
	\begin{itemize}
		\item Um zu testen, ob eine Kurve supersingulär und damit kryptographisch ungeeignet ist, muss die Gruppenordnung $\#E(\FF_{p^r})$ der elliptischen Kurve bestimmt werden -- typischerweise mit dem Schoof-Algorithmus.
		\item Der Begriff supersingulär hat nichts mit singulären Punkten zu tun:
		elliptische Kurven sind per Definition nicht-singulär.
	\end{itemize}
\end{bem}

\begin{bsp}
\label{bsp_19.21}
	Die Kurve $E(\FF_2)\colon y^2 + y = x^3 + x + 1$ ist supersingulär, da $E(\FF_2) = \{ \oh \}$.
\end{bsp}

\begin{bem}
	Supersingularität bleibt bei Übergang zu einem Erweiterungkörper erhalten:
	Ist $E(\FF_{p^r})$ supersingulär, dann auch $E(\FF_{p^{rl}})$ für alle $l \geq 1$. (ohne Beweis)
\end{bem}

\begin{satz}[Erstes Kriterium für Supersingularität]
	Sei $p \geq 3$, $E(\FF_p) \colon y^2 = x^3 + ax^2 + bx + c =: h(x)$ elliptische Kurve.
	Dann ist $E(\FF_{p})$ genau dann supersingulär, wenn der Koeffizient vor $T^{p-1}$ ind $h(T)^{\frac{p-1}{2}} \in \FF_p[T]$ gleich $0$ ist.
\end{satz}

\begin{satz}[Zweites Kriterium für Supersingularität]
	Sei $p = 2$, $E(\FF_{2^r}) \colon y^2 + a_1 xy + y = x^3 + a_2 x^2 + a_4 x + a_6$.
	Dann ist $E(\FF_{2^r})$ genau dann supersingulär, wenn $a_1 = 0$. (ohne Beweis)
\end{satz}

\begin{bsp}
	Siehe Beispiel~\ref{bsp_19.21}, und $E(\FF_p) \colon y^2 = x^3 + x$ ist für $p \kon 3 \modu 4$ supersingulär, denn: \\
	$(T^3 + T)^{(p-1)/2} = \sum_{j=1}^{(p-1)/2} \binom{(p-1)/2}{j} T^{3j} T^{(p-1)/2-j}$ mit $\frac{p-1}{2} + 2j = p-1 \Leftrightarrow 2j = \frac{p-1}{2}$, d.h. wenn $2 \mid \frac{p-1}{2} \Leftrightarrow p \kon 1 \modu 4$ \\
	$\Rightarrow$ Koeffizient vor $T^{p-1}$ ist $\binom{(p-1)/2}{(p-1)/4} \neq 0$ in $\FF_p$. Für $p \kon 3 \modu 4$ kommt $T^{p-1}$ nicht vor, also Koeffizient ist $0$.
\end{bsp}

\begin{bem}
	Der MOV-Algorithmus nutzt bei einer supersingulären Kurve $E(\FF_{p^r})$ aus, dass $t = p^r + 1 - \#E(\FF_{p^r})$ nur einen der Werte $t \in \{0, \pm \sqrt{p^r}, \pm \sqrt{2p^r}, \pm \sqrt{3p^r}, \pm 2 \sqrt{p^r}\}$ annehmen kann.
\end{bem}

\begin{bem}
	Beim \Index{SSSA-Verfahren} (Autoren: Satoh, Smart, Semaev, Araki) handelt es sich um einen schnellen Algorithmus zur Lösung des DL-Problems auf anormalen elliptischen Kurven, welche deswegen kryptographisch ungeeignet sind.
	Die Grundidee ist, die elliptische Kurve über $\FF_p$ als eine über $\QQ_p$ zu betrachten, dem Körper der $p$-adischen Zahlen, und die Logarithmen auf eine Division in $\ZZ_p$ zurückzuführen (was leicht ist).
\end{bem}

\begin{defn}[anormal]
	Eine elliptische Kurve $E(\FF_p)$ heißt \Index{anormal}, wenn $\#E(\FF_p) = p$ ist.
	Dies lässt sich wieder durch Bestimmung von $\#E(\FF_p)$ mit dem Schoof-Algorithmus leicht überprüfen.
	Der SSSA-Algorithmus kann auf Kurven über $\FF_{p^r}$ übertragen werden.
	Er hat polynomielle Laufzeit.
\end{defn}

\subsubsection{Fazit: geeignete elliptische Kurven und Vergleich mit anderen Public-Key-Verfahren}
\begin{bem}
	Eine elliptische Kurve $E(\FF_{p}): y^2 = x^3 + ax + b$ mit vorgegebener Bitzahl für $p$ ist leicht zu finden -- mit Zufallszahlengenerator und Primzahltest, was auch für große Zahlen mit mehreren hundert Dezimalstellen schnell machtbar ist; dafür kennt man ganz gute Algorithmen.
\end{bem}

\begin{bem}
	Man wählt so lange die Parameter $p,a,b$ neu, bis die Diskriminante $4a^3+27b^2$ nicht durch $p$ teilbar ist und somit eine elliptische Kurve vorliegt.
	Ziemlich sicher liegt dann eine kryptographisch geeignete Kurve vor.
	Das testet man nach Berechnen der Gruppenordnung $\#E(\FF_p)$ mit dem Schoof-Algorithmus:
\end{bem}

\begin{bem}
	\begin{itemize}
		\item Ist $\#E(\FF_p)$ glatt, d.h. hat $\#E(\FF_p)$ nur kleine Primteiler, ist die Kurve ungeeignet wegen Silver-Pohlig-Hellman.
		\item Ist $\#E(\FF_p) = p+1$, d.h. die Kurve ist supersingulär, dann ist die Kurve ungeeignet (MOV).
		\item Ist $\#E(\FF_p) = p$, d.h. die Kurve anormal, ist die Kurve ungeeignet (SSSA).
	\end{itemize}
	Ob die Kurve supersingulär oder anormal ist, kann man meist leicht daran erkennen durch Wahl von Punkten $P \in E(\FF_p)$ und dem Test, ob $(p+1)P = \oh$ bzw. $pP = \oh$ gilt.
\end{bem}

\begin{bem}
	Die Wahl eines Punktes $P$ mit nicht zu kleiner Ordnung $n$ muss dann gewährleistet werden.
	Speziell darf $n$ kein Teiler von $p^{rl} - 1$ sein, wenn das DL-Problem in $(\FF_{p^{rl}}^*,\cdot)$ leicht zu lösen ist, und $n$ darf auch kein Vielfaches von $p$ sein (wegen SSSA).
	Auch sollte $n$ nicht glatt sein; man wählt in der Praxis meist Punkte $P$, für die $n = \ord(P)$ eine hinreichend große Primzahl ist; für sie sollte etwa $n > 2^{160}$ gelten.
\end{bem}

\begin{bem}
	Die für allgemeine elliptische Kurven, die in diesem sinne als kryptographisch sicher gelten, bekannten Implementationen des DL-Problems sind alle von exponentieller Komplexität.
	Ein Kryptographieverfahren wie ElGamal bzw. DSA gilt dann als kryptographisch sicher.
\end{bem}

\begin{bem}
	Für konventionelle Kryptoverfahren (RSA und ElGamal/DSA auf $(\FF_{p^r}^*,\cdot)$) gibt es subexponentielle Verfahren zur Lösung des DL-Problems. Dieser Vergleich schlägt sich in der Wahl der Schlüssellängen ($=$ Bitzahl der Größe des endlichen Körpers) nieder:
	Die Schlüssellänge eines elliptischen Kurven-Systems wächst etwas schneller als die dritte Wurzel der Schlüssellänge eines konventionellen Krypto-Systems mit ähnlicher kryptographischer Sicherheit: \todo{Grafik einfügen!}
\end{bem}

\begin{bem}
	\begin{itemize}
		\item Man geht davon aus, dass Kurven $E(\FF_p)$ mit $p \approx 2^{173}$. wo $\#E(\FF_p)$ einen Primteiler $\geq 2^{160}$ hat, die gleiche Sicherheit wie ein RSA-System mit $1024$ Bit bietet (für $4096$ Bit beim RSA nur etwa $313$ bei EC-System!).
		\item Durch die geringere Schlüssellänge bei Verfahren mit elliptischen Kurven kann man diese leicht auf Smart-Cards ohne Koprozessor implementieren.
		Solche Smart-Cards sind wesentlich billiger als Chip-Karten mit Koprozessor.
	\end{itemize}
\end{bem}

\begin{bem}
	Bedenken der elliptischen Kurven-Kryptographie:
	Die Nichteignung supersingulärer und anormaler Kurven kam schnell und überraschend.
	Es ist unklar, ob noch weitere ungeeignete Kurvenfamilien existieren und mit einem schnellen DL-Algorithmus angreifbar sind.
\end{bem}				% Definitionsgebiet
% % % % % % % % % % % 29. Mai
\section{Einschub: Satz von Gauß}
\label{sec:para5}
Wir möchten letztendlich auch nicht-glatte Funktionen $u$ als (verallgemeinerte) \marginnote{29. Apr} Lösungen zulassen. Hierzu nutzen wir Integration, um unerwünschte Ableitungen zu eliminieren. Die Grundidee im Eindimensionalen ist folgende: Gegeben sei:
\[ u'(x) = f(x,u(x)) \text{ auf } \Omega = (a,b) \]
Multipliziere beide Seiten mit einer glatten Funktion $\psi$ und integriere partiell:
\[ u(b)\psi(b) - u(a)\psi(a) - \int\limits_{a}^{b} \psi'u dx = \int\limits_{a}^{b} f(x,u(x)) dx \]
Diese Gleichung macht auch für nicht differenzierbares $u$ Sinn. Im $\RR^n$ benötigen wir den folgenden Satz:

\begin{thm}[Satz von Gauß] \label{thm_13} \label{gauss}
	Sei $\Omega \subseteq \RR^n$ offen mit stückweise stetig differenzierbarem Rand. Ist $u$ stetig differenzierbar, gilt \marginnote{[13]}
	\[ \int_{\Omega} u_{x_i}(x) dx = \int_{\der \Omega} u(x) \nu_i(x) dx \]
	für die nach außen zeigende Einheitsnormale $\nu$ auf $\der\Omega$.
\end{thm}
	
\minisec{Beweis}
	\begin{enumerate}
		\item Lokalisierung: Überdecke $\der \Omega$ durch offene Mengen $U_1,\dots,U_l$, sodass $\der \Omega \cap U_j$ der Graph einer stückweise stetig differenzierbaren Funktion ist und $\Omega \cap U_j$ auf einer Seite des Graphen liegt. Füge $U_0$ mit $\overline{U_0} \subseteq \Omega$ hinzu, sodass $U_0, \dots, U_l$ ganz $\Omega$ überdecken. Wähle hierzu eine Partition der Eins, d.h. unendlich oft differenzierbare Funktionen $\eta_0,\dots,\eta_l$ mit Träger in $U_0, \dots, U_l$ und $\sum_{j=0}^{l} \eta_j = 1$. Es bleibt zu zeigen, dass $\sum_{j=0}^l \int_{\Omega} (\eta_j u)_{x_i} dx = \sum_{j=0}^l \int_{\der \Omega} \eta_j u \nu_i dx$ bzw.
		\[ \int_{\Omega} \nabla(\eta_j u) dx = \int_{\der \Omega} \eta_j u\nu dx \]
		\item Koordinatentransformation: Nach einer orthogonalen Koordinatentransformation dürfen wir annehmen:
		\[ U_j \cap \Omega = \{ x \in \Omega : x_n > g(\tilde{x}) \} \]
		für $\tilde{x} = (x_1,\dots,x_{n-1})$ und eine stückweise glatte Funktion $g \colon \RR^{n-1} \rightarrow \RR$. Die Normale an $\der \Omega \cap U_j$ ist dann gegeben durch
		\[ \nu(\tilde{x},g(\tilde{x})) = \frac{1}{\sqrt{1+|\nabla g(\tilde{x})|^2}} \cdot \begin{pmatrix} \nabla g(\tilde{x}) \\ -1 \end{pmatrix}. \]
		\item Partielle Integration: Mit $v(\tilde{x},h) := (\eta_j u)(\tilde{x},g(\tilde{x})+h)$ ist
		\begin{align}
			&\der_n v(\tilde{x},h) = \der_n (\eta_j u)(\tilde{x},g(\tilde{x}) + h), \notag \\
			&\der_i v(\tilde{x},h) = \der_i (\eta_j u)(\tilde{x},g(\tilde{x}) + h) + \der_i g(\tilde{x}) \der_n(\eta_j u)(\tilde{x},g(\tilde{x}) + h) \notag
		\end{align}
		und somit
		\begin{equation}
		\begin{aligned}
			\int_\Omega \nabla (\eta_j u)(x) dx &= \int_{\RR^{n-1}} \int_{0}^{\infty} \nabla(\eta_j u)(\tilde{x},g(\tilde{x}) + h) dh d\tilde{x} \\ \notag
			&= \int_{\RR^{n-1}} \int_{0}^{\infty} \left(\nabla v - \der_n v \begin{pmatrix} \nabla g \\ 0 \end{pmatrix} \right) (\tilde{x},h) dh d\tilde{x} \\
			&= \sum\limits_{i=1}^{n-1} \left( \int_{0}^{\infty} \int_{\RR^{n-1}} \der_i v(\tilde{x},h) d\tilde{x}, dh \right) e_i - \int_{\RR^{n-1}} \left( \int_{0}^{\infty} \der_n v(\tilde{x},h) dh \right)
		\end{aligned}
		\end{equation}
		Durch partielle Integration bzgl. $x_i$ für $i < n$ ergibt sich aufgrund des kompakten Trägers von $v$:
		\[ \int_{\RR^{n-1}} \der_i v d\tilde{x} = 0 \]
		und durch partielle Integration bzgl. $h$:
		\[ \int_{0}^{\infty} \der_n v(\tilde{x},h) dh = -v(\tilde{x},0) = -(\eta_j u)(\tilde{x},g(\tilde{x})) \] \qed
	\end{enumerate}
	
\begin{bem}[Varianten des Satzes von Gauß]\label{bem_14}
	Der Satz hat verschiedene Varianten: \marginnote{[14]}
	\begin{itemize}
		\item $\int_\Omega \nabla u(x) dx = \int_{\der \Omega} u(x) \nu(x) dx$
		\item $\int \diver b(x) dx = \int_{\der \Omega} b(x) \cdot \nu(x) dx$ für stetig differenzierbares $b\colon \Omega \rightarrow \RR^n$
		\item $\int_\Omega uv_{x_i} + vu_{x_i} dx = \int_{\der \Omega} uv \nu_i dx$ für $u,v$ stetig differenzierbar
		\item $\int_\Omega u \nabla v dx = \int_{\der \Omega} u v \nu dx - \int_\Omega v \nabla u dx \dots$
		\item Der Satz gilt auch für Lipschitz-Rand $\der \Omega$ und schwach differenzierbare Funktionen.
	\end{itemize}
\end{bem}
\newpage				% Satz von Gauß
\section{Schwache Lösungen}
\label{sec:para6}
Frage: Wie können auch nicht-glatte Funktionen $u$ als (verallgemeinerte) Lösungen aufgefasst werden? \marginnote{2. Mai} \\
Idee: Nutze partielle Integration, um Ableitungen loszuwerden.\\
Häufig kann die PDGL in \Index{Erhaltungsform} oder \Index{Divergenzform} geschrieben werden:
\begin{equation}
	\diver G(x,u(x)) = R(x,u(x)) \label{eq_0205_11}
\end{equation}
für Funktionen $G\colon \RR^n \times \RR \rightarrow \RR^n, R\colon \RR^n \times \RR \rightarrow \RR$. Wir wollen \eqref{eq_0205_11} auf einem Gebiet $\Omega$ lösen für Cauchy-Daten
\[u = g \text{ auf } \Gamma \subseteq \der \Omega \]

\begin{bsp} \label{bsp_15}
	$a(x)u_{x_1} + u u_{x_2} = c(x,u)$ ist äquivalent \marginnote{[15]} zu $(au)_{x_1} + \left( \frac{u^2}{2} \right)_{x_2} = c(x,u) + a_{x_1} u$ und kann in Erhaltungsform geschrieben werden für $G(x,u) = \begin{pmatrix} au \\ u^2/2 \end{pmatrix}, R(x,u) = c(x,u) + a_{x_1} u$. \\
	Wir multiplizieren \eqref{eq_0205_11} mit einer sogenannten \Index{Testfunktion}
	\[ \psi \colon \RR^n \rightarrow \RR, \psi \text{ ist glatt}, \psi = 0 \text{ auf } \der \Omega \setminus \Gamma \]
	und integrieren partiell:
	\begin{equation}
	\begin{aligned}
		&\int_\Omega \psi(x) (\diver G(x,u(x)) - R(x,u(x))) dx \\ \notag
		= &\int_{\der \Omega} \psi(x) G(x,u(x)) \cdot v(x) dx - \int_\Omega \nabla \psi(x) \cdot F(x,u(x)) + \psi(x) R(x,u(x)) dx \\
		= &\int_\Gamma \psi(x) G(x,g(x)) \cdot v(x) dx - \int_\Omega \nabla \psi(x) \cdot G(x,u(x)) + \psi(x) R(x,u(x)) dx
	\end{aligned}
	\end{equation}
	Die Gleichung macht auch für nicht-differenzierbares $u$ Sinn.
\end{bsp}

% % % % % % % % % % % % % 02. Mai	
\begin{defn}[schwache Lösung]
\label{def_schwache_lsg} \label{def_16}
	Eine messbare, beschränkte Funktion $u(x)$ auf $\Omega$ heißt schwache Lösung von \eqref{eq_0205_11}, wenn sie \marginnote{[16]}
	\begin{equation}
		0 = \int_\Gamma \psi(x) G(x,g(x)) \cdot v(x) dx - \int_\Gamma \nabla \psi(x) \cdot G(x,u(x)) + \psi (x) R(x,u(x)) dx \label{eq_0205_12}
	\end{equation}
	für alle Testfunktionen $\psi$ erfüllt.
\end{defn}

\minisec{Unstetige Lösungen und Schocks}
	Seien
	\begin{itemize}
		\item $G \subseteq \Omega$ stückweise glatte Hyperfläche, die $\Omega$ in $\Omega_1$ und $\Omega_2$ teilt.
		\item $u$ eine schwache Lösung von \eqref{eq_0205_11}, die auf $\Omega_1,\Omega_2$ stetig differenzierbar ist
		\item $\nu_i$ die äußere Einheitsnormale an $\der \Omega_i$,$i=1,2$
		\item $[f]_-^+$: Sprung einer Funktion $f$ über $C$ von $\Omega_1$ nach $\Omega_2$
	\end{itemize}
	Für alle Testfunktionen $\psi$ gilt:
	\begin{equation}
	\begin{aligned}
		0 &= \int_\Gamma \psi G(x,g(x)) \cdot \nu(x)dx - \int_{\Omega_1} \nabla \psi \cdot G(x,u(x)) + \psi R(x,u(x)) dx - \int_{\Omega_2} \nabla \psi \cdot G(x,u(x)) + \psi R(x,u(x)) dx \\ \notag
		&= \int_{\Gamma} \psi G(x,g(x)) \cdot \nu(x) dx - \int_{\der\Omega_1} \psi G(x,u(x)) \cdot \nu_1(x) dx + \int_{\Omega_1} \psi (\diver G(x,u(x)) - R(x,u(x))) dx \\
		&\quad - \int_{\der\Omega_2} \psi G(x,u(x)) \cdot \nu_2(x) dx + \int_{\Omega_2} \psi (\diver G(x,u(x)) - R(x,u(x))) dx \\
		&= \int_C \psi[G(x,u(x))]_-^+ \cdot \nu_1(x) dx,
	\end{aligned}
	\end{equation}
	da $\diver G(x,u(x)) - R(x,u(x)) = 0$ in $\Omega_1,\Omega_2$. Wir erhalten
	\begin{equation}
		[G(x,u(x))]_-^+ \cdot \nu(x) = 0 \label{eq_0205_13}
	\end{equation}
	für die Normale $\nu$ an $C$. Eine Unstetigkeit von $u$ bzw. $G$, ein sogenannter \Index{Schock}, kann also nur entlang einer Fläche $C$ normal zu $[G]_-^+$ auftreten. \\
	\mbox{} \\
	In 2D kann die Kurve $C$ lokal als Graph einer Funktion $x_2(x_1)$ (bzw. $x_1(x_2)$) geschrieben werden. Ein Tangentialvektor ist somit $\begin{pmatrix} 1 \\ dx_2/dx_1 \end{pmatrix}$, ein Normalvektor $\begin{pmatrix} dx_2/dx_1	\end{pmatrix}$, und \eqref{eq_0205_13} wird zur sogenannten \bet{Rankine-Hugoniot} \bet{-Bedingung} für einen Schock: \index{Rankine-Hugoniot-Bedingung}
	\begin{equation}
		\frac{dx_2}{dx_1} = \frac{[G_2]_-^+}{[G_1]_-^+} \label{eq_0205_14}
	\end{equation}
	
\begin{bem}[Rankine-Hugoniot-Bedingung für semilineare PDGL] \label{bem_17}
	Für semilineare Gleichungen \marginnote{[17]} $0 = b(x) \cdot \nabla u(x) + c(u(x),x)$ wird dies zu $\frac{dx_2}{dx_1} = \frac{b_2[u]_-^+}{b_1[u]_-^+} = \frac{b_1}{b_2}$, d.h. der Schock liegt entlang einer charakteristischen Kurve.
\end{bem}
	
\begin{bsp} \label{bsp_18}
	Betrachte das Cauchy-Problem \marginnote{[18]}
	\[ \begin{cases}
		u_t+uu_x = 0 \text{ auf } (0,\infty) \times \RR \\
		u(0,x) = g(x) := \begin{cases}
			1	& \text{ für } x \leq 0 \\
			1-x & \text{ für } 0 \leq x \leq 1 \\
			0	& \text{ für } 1 \leq x \end{cases} \end{cases} \]
	Entlang der charakteristischen Kurven $s \mapsto (s,y+g(y)s)$ für $y \in \RR$ ist $u = z(s) = g(y)$ konstant, somit
	\[ u(t,x) = \begin{cases}
		1	&	\text{ für } x \leq t \\
		\frac{1-x}{1-t} & \text{ für } t \leq x \leq 1 \\
		0	&	\text{ für } 1 \leq x \end{cases} \]
	solange $t < 1$. In $(t,x) = (1,1)$ schneiden sich die charakteristischen Kurven, und die Lösung mithilfe der Charpits Gleichungen bricht zusammen. An dieser Stelle setzen wir einen Schock entlang einer Kurve $x = X(t)$ ein und betrachten die Erhaltungsform $u_t + \left( \frac{u^2}{2} \right)_x = 0$. Sei $(0,y_\pm(t))$ der Anfangspunkt der charakteristischen Kurven, die $X(t)$ zur Zeit $t$ von links ($-$) und von rehts ($+$) schneiden. Es gilt somit
	\[ X(t) = y_-(t) + g(y_-(t))t = y_+(t) + g(y_+(t))t \]
	Weiterhin gilt die Rankine-Hugoniot-Bedingung
	\[ \frac{dX}{dt} = \frac{\frac{u_+^2}{2}-\frac{u_-^2}{2}}{u_+ - u_-} = \frac{u_+ + u_-}{2} = \frac{g(y_+(t)) + g(y_-(t))}{2} \]
	und $X(1)=1$. Eine Lösung dieser Gleichungen ist gegeben durch $X(t) = \frac{1+t}{2}, y_-(t) = \frac{1-t}{2}, y_+(t) = \frac{1+t}{2}$ (die Charakteristiken von links und rechts haben $u_-=1, u_+ = 0$), folglich ist
	\[ u(t,x) = \left. \begin{cases}
		1	&	\text{ für } x \leq t \\
		\frac{1-x}{1-t}	&	\text{ für } t \leq x \leq 1 \\
		0	&	\text{ für } 1 \leq x \end{cases} \right\} \text{ für } t<1 \text{ und } u(t,x) = \left. \begin{cases}
		1	&	\text{ für } x \leq \frac{1+t}{2} \\
		0	&	\text{ für } \frac{1+t}{2} \leq x \end{cases} \right\} \text{ für } t \geq 1 \]
	eine schwache Lösung des Problems.
\end{bsp}
	
\begin{bem} \label{bem_19}
	Schwache Lösungen sind im Allgemeinen nicht eindeutig. Betrachte z.B. \marginnote{[19]}
	\[ \begin{cases} 
		u_t + uu_x = 0 \text{ auf } (0,\infty) \times \RR \\
		u(0,x) = g(x) := \begin{cases}
			0 & \text{ für } x \leq 0 \\
			1 & \text{ sonst} \end{cases} \end{cases} \]
	mit den schwachen Lösungen:
	\begin{itemize}
		\item $u(t,x) = \begin{cases}
			0 & \text{ für } x \leq t/2 \\
			1 & \text{ sonst} \end{cases}$ \\
		(Schock zu Erhaltungsform $u_t + \left( \frac{u^2}{2} \right)_x = 0$; Rankine-Hugoniot: $\frac{dx}{dt} = \frac{\frac{u_+^2}{2}-\frac{u_-^2}{2}}{u_+-u_-} = \frac{u_++u_-}{2}$)
		\item $u(t,x) = \begin{cases}
			0 & \text{ für } x \leq 2t/3 \\
			1 & \text{ sonst} \end{cases}$ \\
		(Schock zu Erhaltungsform $\left( \frac{u^2}{2} \right)_t + \left( \frac{u^3}{2} \right)_x = 0$; Rankine-Hugoniot: $\frac{dx}{dt} = \frac{2}{3} \frac{u_+^3-u_-^3}{u_+^2-u_-^2} = \frac{2}{3} \frac{u_+^2+u_-u_++u_-^2}{u_++u_-}$)
		\item $u(t,x) = \begin{cases}
			0 & \text{ für } x\leq 0 \\
			x/t & \text{ für } 0 < x < t \\
			1 & \text{ sonst} \end{cases}$ \\
		(Verdünnungswelle)
	\end{itemize}
	Welche Erhaltungsform korrekt ist, hängt vom modellierten physikalischen Problem ab -- auch, ob ein Schock oder stattdessen eine Verdünnungswelle korrekt sind. Man kann für eine vorgegebene Erhaltungsform eine eindeutige Lösung selektieren durch Hinzunahme einer zusätzlichen Bedingung, zum Beispiel: \begin{enumerate}
		\item Entropiebedingung: Wir legen eine bestimmte Funnktion von $f(u)$ fest, die über den Schock zunehmen soll, d.h. $[f(u)]_-^+ \geq 0$. Die Wahl $f(u) = u$ verbietet obige Schocks.
		\item Viskosität: Wir betrachten $0 = F(\nabla u, u, x)$ als Grenzfall für $\varepsilon \rightarrow 0$ von $\varepsilon \Delta u = F(\nabla u,u,x)$, welches eine eindeutige Lösung besitzt.
		\item Kausalität: Ist eine Variable die Zeit, so fließt physikalisch Information entlang der charakteristischen Kurven in Richtung wachsender Zeit. Es soll Information in den Schock hinein und nicht aus ihm hinaus fließen (dies verbietet obige Schocks).
	\end{enumerate}
\end{bem}

% % % % % % % % % % % % 6. Mai	
\minisec{Viskositätslösungen}
	Betrachte die nichtlineare PDGL 1. Ordnung \marginnote{6. Mai}
	\begin{equation}
		H(x,u(x),\nabla u(x)) = 0 \text{ auf } \Omega \label{eq_0605_15}
	\end{equation}
	mit Randbedingungen auf $\der \Omega$, wobei $H\colon \Omega \times \RR \times \RR^n \rightarrow R$ stetig ist und konvex in $\nabla u$. \eqref{eq_0605_15} heißt auch \Index{Hamilton-Jacobi-Gleichung}. Lösungen sind typischerweise weder eindeutig, noch überall differenzierbar.
	
\begin{bsp}[Eikonalgleichung] \label{bsp_20}
	Die \Index{Eikonalgleichung} \marginnote{[20]}
		\[ 0 = H(x,u,p) = |p|-1 \]
	wird oft verwendet, um den Abstand $u(x) = \dist(x,\Gamma)$ einer Teilmenge $\Gamma \subset \RR^n$ zu berechnen. Allerdings existiert bereits für $n = 1$ und Randdaten $u = 0$ auf $\Gamma = \{0,1\}$ keine stetig differenzierbare Lösung $u$, aber viele Funktionen $u$, die \eqref{eq_0605_15} fast überall erfüllen. \\
	\mbox{} \\
	Idee, um Wohlgestelltheit zu erhalten: Füge eine \glqq Viskositätsterm\grqq \ $-\varepsilon \Delta u_\varepsilon$ hinzu und betrachte $\varepsilon \rightarrow 0$:
	\begin{equation}
		-\varepsilon \Delta u_\varepsilon + H(x,u_\varepsilon(x),\nabla u_\varepsilon (x)) = 0 \label{eq_0605_16}
	\end{equation}
\end{bsp}
	
\begin{lemma} \label{lemma_21}
	Sei $\Omega \subset \RR^n$ offen, \marginnote{[21]} $\overline{x} \in \Omega$, $v, v_k \in C^1(\Omega)$ mit $v_k \rightarrow v$ lokal gleichmäßig. Hat $v$ ein striktes lokales Minimum in $\overline{x}$, dann existiert eine Folge $x_k \rightarrow \overline{x}$ derart, dass $v_k$ ein lokales Minimum in $x_k$ hat.
\end{lemma}
	
\minisec{Beweis}
	Sei o.B.d.A. $\overline{x} = 0, v(\overline{x})=0$. \begin{itemize}
		\item Definiere $\omega_\rho := \min_{\der B_\rho(0)} v > 0$ für $0 < \rho < r$.
		\item Definiere $\rho_k := \inf\{ \rho : \frac{\omega_\rho}{4} \geq \sup_{B_r(0)} |v_k-v|\}$, dann gilt $\rho_k \rightarrow 0$.
	\end{itemize}
	Angenommen, es wäre $\rho_{k_j} \geq \delta > 0$ für eine Folge $k_j \rightarrow \infty$, dann wäre $\frac{\omega_\delta}{4} < \sup_{B_r(0)} |v_{k_j} - v| \rightarrow 0$. Widerspruch! \\
	1. Fall: $\rho_k = 0$. \\
	$\Rightarrow v_k = v \Rightarrow x_k = \overline{x}$ ist lokale Minimalstelle von $v_k$. \\
	2. Fall: $\rho_k > 0$. \\
	$\Rightarrow \min_{\der B_{\rho_k}(0)} v_k \geq \min_{\der B_{\rho_k}(0)} v - \frac{\omega_{\rho_k}}{4} = \frac{3 \omega_{\rho_k}}{4} \geq \frac{\omega_{\rho_k}}{4} \geq v_k(0) \Rightarrow v_k$ hat ein $x_k \in B_{\rho_k}(0)$ als Minimalstelle. \\
	Damit folgt $x_k \rightarrow 0$. \qed

\mbox{} \\
Sei $u_\varepsilon \in C^2(\Omega)$ eine Lösung von \eqref{eq_0605_15} mit $u_\varepsilon \rightarrow u$ lokal gleichmäßig. Betrachte eine Vergleichsfunktion $\phi \in C^2(\Omega)$ mit $\phi \leq u$, sodass $u-\phi$ ein striktes Minimum in $\overline{x} \in \Omega$. \\
Lemma \ref{lemma_21} impliziert, dass eine Folge $\varepsilon_k \rightarrow 0$, $x_k \rightarrow \overline{x}$ existiert, sodass $u_{\varepsilon_k} - \phi$ ein lokales Minimum in $x_k$ besitzt, das heißt:
\[ \begin{array}{rcl}
	\nabla(u_{\varepsilon_k}(x_k) - \phi(x_k)) = 0	&	\hspace*{1cm}	&	\text{(Ableitung = 0)} \\
	D^2(u_{\varepsilon_k}(x_k) - \phi(x_k)) \geq 0	&	&	\text{(Hesse-Matrix positiv semidefinit)} \\
	\Delta u_{\varepsilon_k}(x_k) - \Delta \phi(x_k) \geq 0	&	& \text{(s.o.)}
\end{array} \]
Damit folgt
\begin{equation}
\begin{aligned}
	0 &= -\varepsilon_k \Delta u_{\varepsilon_k}(x_k) + H(x_k,u_{\varepsilon_k}(x_k)) \\ \label{eq_0605_17}
	&\leq -\varepsilon_k \Delta\phi (x_k) + H(x_k,u_{\varepsilon_k}(x_k),\nabla \phi (x_k)) \rightarrow H(\overline{x},u_{\varepsilon_k}(\overline{x}),\nabla \phi (\overline{x}))
\end{aligned}
\end{equation}
das heißt
\[ H(\overline{x},u_{\varepsilon_k}(\overline{x})) \geq 0 \text{, falls } u-\phi \text{ ein striktes lokales Minimum in } \overline{x} \text{ besitzt.} \]
Dies motiviert zu folgendem Begriff:

\begin{defn}[Viskositätslösung] \label{def_22}
	Sei $\Omega \subset \RR^n$ offen, $H\colon \Omega \times \RR \times \RR^n$ stetig. \marginnote{[22]}
	\begin{enumerate}
		\item $u \in C^0(\Omega)$ heißt \bet{Viskositäts-Superlösung}, falls $H(x,u(x),\nabla\phi(x)) \geq 0$ für alle $\phi \in C^1(\Omega)$ derart, dass $u-\phi$ ein Minimum in $x$ besitzt.
		\item $u \in C^0(\Omega)$ heißt \bet{Viskositäts-Sublösung}, falls $H(x,u(x),\nabla\phi(x)) \leq 0$ für alle $\phi \in C^1(\Omega)$ derart, dass $u-\phi$ ein Maximum in $x$ besitzt.
		\item $u \in C^0(\Omega)$ heißt \Index{Viskositätslösung}, falls $u$ Viskositäts-Superlösung und Viskositäts-Sublösung ist.
	\end{enumerate}
\end{defn}

\begin{bsp}[Abstandsfunktion]\label{lemma_23}
	Für $\Gamma \subset \RR^n$ abgeschlossen ist $u(x) = \dist(x,\Gamma)$ eine Viskositätslösung von $|\nabla u| - 1 = 0$ auf $\RR^n \setminus \Gamma$. \marginnote{[23]}
\end{bsp}
	
\minisec{Beweis}
	\begin{enumerate}
		\item zur Sublösung: Sei $x$ ein lokales Maximum von $u - \phi$, das heißt
		\begin{equation}
		\begin{aligned}
			&(u-\phi)(x+z)-(u-\phi)(x) \leq 0 \\ \notag
			\Leftrightarrow \quad &\phi(x+z) - \phi(x) \geq u(x+z) - u(x) \geq -|z| \\
			\Rightarrow \quad &\frac{\phi(x+s\zeta) - \phi(x)}{s} \geq -|\zeta| \text{ und somit } \nabla\phi(x) \cdot \zeta \geq -|\zeta| \ \forall \zeta \in \RR^n 
		\end{aligned}
		\end{equation}
		Wähle $\zeta = - \frac{\nabla \phi(x)}{|\nabla \phi(x)|}$ und erhalte $|\nabla \phi(x)| \leq 1$ oder $|\nabla \phi(x)| - 1 \leq 0$.
		\item zur Superlösung:  Sei $x$ ein lokales Minimum von $u - \phi$, das heißt
		\begin{equation}
		\begin{aligned}
			&(u-\phi)(x+z)-(u-\phi)(x) \geq 0 \\ \notag
			\Leftrightarrow \quad &\phi(x+z) - \phi(x) \leq u(x+z) - u(x) \\
		\end{aligned}
		\end{equation}
		Sei $\overline{x} \in \Gamma$ mit $u(x) = |x-\overline{x}|$ und wähle $z = -s \frac{x-\overline{x}}{|x-\overline{x}}$.
		\begin{equation}
		\begin{aligned}
			&\Rightarrow \quad \frac{ \phi \left(x+s\frac{\overline{x}-x}{|\overline{x}-x|} \right) - \phi(x)}{s} \leq \frac{|\overline{x} - x| - s - |\overline{x}-x|}{s} = -1 \\ \notag
			&\Rightarrow \quad -\nabla \phi(x) \cdot \frac{\overline{x}-x}{|\overline{x}-x|} \geq 1
			\quad \Rightarrow \quad |\nabla\phi(x)| \geq 1 \text{ oder } |\nabla \phi(x)|-1 \geq 0. 
		\end{aligned}
		\end{equation}
		\qed
	\end{enumerate}
	
\begin{bem} \label{bem_24}
	Im Allgemeinen ist $u  = -\dist(x,\Gamma)$ \marginnote{[24]} keine Viskositätslösung von $|\nabla u|-1 = 0$: Sei $u_{\min} = u(\overline{x})$ ein lokales Minimum von $-\dist(x,\Gamma)$ und $\phi \equiv u_{\min}$, sodass $u  -\phi$ ein lokales Minimum in $\overline{x}$ hat. Dann ist $|\nabla \phi| - 1 = -1 < 0$ und $u$ damit keine Viskositäts-Superlösung. \\
	Allerdings ist $u$ eine Viskositätslösung von $1-|\nabla u| = 0$.
\end{bem}
	
	\mbox{} \\
	Wir betrachten nun Hamilton-Jacobi-Gleichungen der Gestalt
	\begin{equation}
		H(x,\nabla u(x)) = 0 \label{eq_0605_18} 
	\end{equation}
	mit $H$ stetig, $H(x,\cdot)\colon \RR^n \rightarrow \RR$ gleichmäßig konvex, $H(x,p) \xlongrightarrow{|p| \rightarrow \infty} \infty$ gleichmäßig in $x$ (Koerzivität) und $H(x,0) \leq 0$.
	
\begin{defn}[Stützfunktion]
	Eine Funktion \marginnote{[25]}
	\begin{equation}
	\begin{aligned}
		L_x\colon \RR^n &\longrightarrow \RR \\ \notag
		w &\longmapsto \sup_{H(x,p) \leq 0} w \cdot p
	\end{aligned}
	\end{equation}
	heißt \Index{Stützfunktion} der Menge $\{p \in \RR^n : H(x,p) \leq 0\}$.
\end{defn}
	
\begin{bsp} \label{bsp_26}
	$H(x,p) = |p| - \frac{1}{v(x)}$ für eine Geschwindigkeit $v(x)$ hat die Stützfunktion $L_x(w) = \frac{|w|}{v(x)}$. \marginnote{[26]}
\end{bsp}
	
\begin{defn}[optische Distanz] \label{def_27}
	\[ \delta(x,y) = \inf \left\{ \int_{0}^{1} L_{c(t)}(\dot{c}(t)) dt : c\colon [0,1] \rightarrow \RR^n, c(0) = x, c(1) = y\right\} \]
	heißt die \Index{optische D<istanz} zwischen $x$ und $y$. \marginnote{[27]}
\end{defn}
	
\begin{bsp} \label{bsp_28}
	Die optische Distanz von $H(x,p) = |p| - \frac{1}{v(x)}$ ist gegeben durch
	\[ \delta(x,y) = \inf_{\substack{c\colon[0,1] \rightarrow \RR^n \\ c(0) = x, c(1) = y}} \int_{0}^{1} \frac{|\dot{c}(t)|}{v(c(t))} dt \]
	und besc<<hreibt die Zeit, die eine seismische Welle mit der Ausbreitungsgeschwindigkeit $v(x)$ benötigt, um von $x$ nach $y$ zu gelangen. \marginnote{[28]}
\end{bsp}

% % % % % % % 09.05.
\begin{thm}[Hopf-Lax-Formel] \label{thm_29}
	Sei $\Omega \subset \RR^n$ offen und beschränkt \marginnote{9. Mai \\ \ [29]} und $g\colon \der \Omega \rightarrow \RR$ mit der Eigenschaft $g(x)-g(y) \leq \delta(y,x)$ für alle $x,y \in \der \Omega$. Dann ist
	\[ u(x) = \inf_{y \in \der \Omega} \{g(y) + \delta(y,x) \} \]
	eine lipschitzstetige Viskositätslösung von $H(x,\nabla u(x)) = 0$ auf $\Omega$ mit $u = g$ auf $\der \Omega$.
\end{thm}
	
\begin{bem} \label{bem_30}
	\begin{itemize}
		\item Dieser Satz impliziert die Existenz einer Viskositätslösung. \marginnote{[30]}
		\item Für $H(x,p) = |p| - \frac{1}{v(x)}$ und $g \equiv 0$ ist $u(x)$ die Ankunftszeit einer seismischen Welle, die auf $\der \Omega$ startet.
		\item $g(x) - g(y) \leq \delta(y,x)$ bedeutet, dass die Wellenfront zum Erreichen von $x$ nicht länger braucht als um von $y$ nach $x$ zu gelangen.
	\end{itemize}
\end{bem}
	
\begin{lemma} \label{bem_31}
	\begin{enumerate}[(i)]
		\item $L_x$ ist konvex. \marginnote{[31]}
		\item $L_x(w) \leq C|w|$ für eine Konstante $C$ unabhängig von $w$.
		\item $L_x(\lambda w) = \lambda L_x(w)$ für alle $\lambda > 0$.
	\end{enumerate}
\end{lemma}
	
\minisec{Beweis}
	\begin{enumerate}[(i)]
		\item $L_x(tq+(1-t)w) = \sup_{H(x,p) \leq 0} [tq + (1-t)w] \cdot p = \sup_{H(x,p) \leq 0} tq \cdot p + (1-t)w \cdot p$ \\
		$\leq t \sup_{H(x,p) \leq 0} q \cdot p + (1-t) \sup_{H(x,p) \leq 0} w \cdot p = tL_x(q) + (1-t)L_x(w)$ für alle $t \in [0,1]$.
		\item $L_x(w) = \sup_{H(x,p) \leq 0} w \cdot p \leq \sup_{H(x,p) \leq 0} |w| \cdot |p| \leq (\sup_{H(x,p) \leq 0} |p|) |w|$
		\item $L_x(\lambda w) = \sup_{H(x,p) \leq 0} \lambda w \cdot p = \lambda \sup_{H(x,p) \leq 0} w \cdot p$ \qed
	\end{enumerate}
	
\begin{bem} \label{bem_32}
	Das obige Lemma impliziert, dass $\delta$ eine Pseudometrik ist, falls $H(x,p) = H(x,-p)$. \marginnote{[32]}
\end{bem}
	
\minisec{Beweis von Theorem \ref{thm_29}}
	Wir betrachten den Fall, dass $H(x,P) \equiv H(p)$ und $H(p) = H(-p)$. In diesem Fall gilt $\delta(x,y) = L(y-x)$, da $L$ nicht von $x$ abhängig ist. \begin{enumerate}[a)]
		\item $u$ ist lipschitzstetig und $u(x) \leq u(y) + \delta(x,y)$ für alle $x,y \in \overline{\Omega}$:
		\begin{equation}
		\begin{aligned}
			u(x) - u(y) &= \inf_{z_1 \in \der\Omega} (g(z_1) + \delta(x,z_1)) - \inf_{z_2 \in \der \Omega} (g(z_2) + \delta(y,z_2)) \\ \notag
			&\leq \sup_{z_2 \in \der \Omega} g(z_2) + \der(x,z_2) - g(z_2) - \der(y,z_2) \\
			&\leq \der(x,y) = L(y-x) \overset{(ii)}{\leq} C|y-x|
		\end{aligned}
		\end{equation}
		\item Sublösung: \\
		$u-\phi$ habe ein lokales Maximum in $x$, d.h. $u(x')-\phi(x') \leq u(x) - \phi(x)$, bzw.
		\[ \phi(x') - \phi(x) \geq u(x') - u(x) \geq -\delta(x,x') \]
		Mit $\zeta$ beliebig und $x' = x + s \zeta$ folgt
		\[ \frac{\phi(x+s\zeta)-\phi(x)}{s} \geq \frac{-\delta(x,x+s\zeta)}{s} \overset{\substack{L \text{ unabh.} \\ \text{von } x}}{=} \frac{-L(s\zeta)}{s} \overset{(iii)}{=} -L(\zeta) \quad \Rightarrow \quad \nabla \phi \cdot \zeta \geq -L(\zeta) \]
		Analog folgt $\nabla \phi \cdot \zeta \leq L(\zeta) = \sup_{H(p) \leq 0} \zeta \cdot p$, also $|\nabla \phi \cdot \zeta| \leq L(\zeta)$ für alle $\zeta \in \RR^n$. Angenommen, $H(\nabla \phi) > 0$. Da $H$ konvex und stetig ist, existiert eine Hyperfläche mit Normale $\nu$, die $\{p : H(p) \leq 0\}$ und $\nabla \phi$ trennt. Damit wäre $|\nabla \phi \cdot \nu| > \sup_{H(p) \leq 0} p \cdot \nu = L(\nu)$. Widerspruch.
		\item Superlösung: \\
		Sei $u(x) = g(\overline{y}) + \delta(x,\overline{y}) = g(\overline{y}) + L(x-\overline{y})$ für ein $\overline{y} \in \der \Omega$ und definiere $c(t) = \overline{y} + t(x-\overline{y})$. Wir haben
		\begin{equation}
		\begin{aligned}
			g(\overline{y}) + L(x-\overline{y}) &= u(x) \leq u(c(t)) + \delta(c(t),x) ) u(c(t)) + (1-t)L(x-\overline{y}) \\ \notag
			&\leq g(\overline{y} + \delta(\overline{y},c(t)) + (1-t)L(x-\overline{y}) = g(\overline{y}) + L(x-\overline{y})
		\end{aligned}
		\end{equation}
		sodass $u(c(t)) = g(\overline{y}) + tL(x-\overline{y}) = g(\overline{y}) + L(c(t)-\overline{y})$. Nun besitze $u-\phi$ ein lokales Minimum in $x$, d.h. $\phi(x) - \phi(x') \geq u(x) - u(x')$, und setze $x=c(1), x'=c(1-s)$. Wir erhalten:
		\begin{equation}
		\begin{aligned}
			&\frac{\phi(x)-\phi(c(1-s))}{s} \geq \frac{u(x) - u(c(1-s))}{s} \\ \notag
			\Rightarrow \quad &\nabla \phi(x) \cdot (x-\overline{y}) = \nabla \phi(x) \cdot \dot{c}(1) \geq \nabla(u) \cdot \dot{c}(1) = L(x-\overline{y}) = \sup_{H(p) \leq 0} (x-\overline{y}) \cdot p
		\end{aligned}
		\end{equation}
		Daher folgt $(\nabla \phi(x) + \alpha(x-\overline{y})) \cdot (x-\overline{y}) > \sup_{H(p) \leq 0} (x-\overline{y})\cdot p$ für alle $\alpha > 0$ und somit $H(\nabla \phi(x) + \alpha(x- \overline{y})) > 0$ und $H(\nabla \phi) \geq 0$ wegen Stetigkeit.
		\item Randdaten: \\
		Sei $x \in \der\Omega$, dann folgt aus $g(x)-g(y) \leq \delta(y,x)$, dass $g(x) \leq g(y) + \delta(y,x)$ für alle $y \in \der\Omega$ und somit $g(x) \leq u(x)$. Des Weiteren ist $u(x) \leq g(x) + \delta(x,x) = g(x)$ und damit $u(x) = g(x)$. \qed
	\end{enumerate}

\mbox{} \\
Für Eindeutigkeit müssen wir zusätzliche Bedingungen an $H$ stellen, wie das folgende Beispiel zeigt. Anschließend formulieren wir zwei mögliche Resultate für die Eindeutigkeit.

\begin{bsp} \label{bsp_33}
	Sei $\Psi \in C^1(\Omega)$, $\Psi = 0$ auf $\der\Omega$,\marginnote{[33]} und betrachte $H(x,p) = |p|^2 - |\nabla \Psi(x)|^2$. Dann ist sowohl $u = \Psi$ als auch $u = -\Psi$ eine (Viskositäts-)Lösung von $0 = H(x,\nabla u (x))$.
\end{bsp}
	
\begin{thm}[Eindeutigkeit durch Vergleich] \label{thm_34}
	Sei $\Omega \subset \RR^n$ beschränkt \marginnote{[34]} und offen sowie $H\colon \overline{\Omega} \times \RR \times \RR^n$ stetig mit \begin{itemize}
		\item $H(x,u,p) - H(x,v,p) > \gamma(u-v)$ für ein $\gamma > 0$.
		\item $|H(x,u,p) - H(y,u,p)| \leq C|y-x|(1+|p|)$ für ein $C > 0$.
	\end{itemize}
	Falls $u$ eine Viskositäts-Sublösung und $v$ eine Viskositäts-Superlösung von $0 = H(x,u(x),\nabla u(x))$ ist mit $u \leq v$ auf $\der \Omega$, dann ist auch $u \leq v$ auf $\Omega$. Damit folgt, dass die Viskositätslösung eindeutig ist.
\end{thm}
	
\minisec{Beweisidee}
	Angenommen, $u$ und $v$ seien glatt und $u-v$ habe ein Maximum in $x_0 \in \overline{\Omega}$ mit $u(x_0) - v(x_0) > 0$. Nach Definition \ref{def_22} haben wir
	\[ \begin{array}{c}
		H(x_0,u(x_0),\nabla v(x_0)) \leq 0 \\
		H(x_0,v(x_0),\nabla u(x_0)) \geq 0
	\end{array}\]
	Wegen $\nabla(u-v)(x_0) = 0$ gilt $\nabla u(x_0) = \nabla v(x_0)$, und damit den Widerspruch
	\[ 0 \geq H(x_0,u(x_0),\nabla v(x_0)) = H(x_0,u(x_0),\nabla(u(x_0)) > H(x_0,v(x_0),\nabla u(x_0)) \geq 0 \]
	Für nichtglatte $u,v$ ist der Beweis aufwändiger. \qed
	
\begin{thm}[Eindeutigkeit für Hamilton-Jacobi-Bellman-Gleichung] \label{thm_35}
	Sei \marginnote{[35]}
	\begin{equation}
	\begin{aligned}
		H\colon (\RR \times \RR^{n-1}) \times \RR^n &\longrightarrow \RR \\ \notag
		((t,x),(p^t,p^x)) &\longmapsto p^t + \tilde{H}(x,p^x)
	\end{aligned}
	\end{equation}
	mit $\tilde{H}$ stetig und \begin{itemize}
		\item $|\tilde{H}(x,p) - \tilde{H}(x-q)| \leq C|p-q|$ für ein $C > 0$
		\item $|\tilde{H}(x,p) - \tilde{H}(y,p)| \leq C|y-x|(1+|p|)$
	\end{itemize}
	Dann existiert höchstens eine Viskositätslösung von $0 = H((t,x),(u_t(t,x),\nabla u(t,x))) = u_t + \tilde{H}(x,\nabla u)$ mit gegebenen Randdaten auf $t = 0$.
\end{thm}

\minisec{Beweis}
	siehe z.B. Evans, \glqq PDEs\grqq, S. 587 
\newpage				% Schwache Lösungen
\section{Fermatsche und Mersennesche Primzahlen}
\label{sec:para7}
	
\minisec{Bemerkung}
	$2^k + 1$ prim $\Rightarrow$ $k$ ist Potenz von $2$ (vgl. Aufgabe 4). \marginnote{19.12. \\ \ [20]}
	
Für $n \in \NN$ sei
\[ F_n := 2^{2^n} + 1 \]
die $n$-te Fermatsche Zahl.

\begin{enumerate}[(1)]
	\item $F_0 = 3, F_1 = 5, F_2 = 17, F_3 = 257$ und $F_4 = 65537$ sind \bet{Fermatsche Primzahlen}. \index{Fermatsche Primzahl}
\end{enumerate}

\minisec{Satz}
	Ein regelmäßiges $n$-Eck ist mit Zirkel und Lineal genau dann konstruierbar, wenn $n$ von der Gestalt $n = 2^\nu p_1 p_2 \cdots p_r$ mit $r \in \NN_0$, paarweise verschiedenen Fermatschen Primzahlen $p_1,\dots,p_r$ und beliebigem $\nu \in \NN_0$ ist.
	
\begin{enumerate}[(1)] \setcounter{enumi}{1}
	\item $F_5$ ist keine Primzahl, sondern durch $641$ teilbar.
	\item Für $n \geq 3$ hat jeder Primteiler von $F_n$ die Gestalt $p = t 2^{n+2} + 1$.
	\item Auch $F_6$ ist nicht prim, sondern teilbar durch $274177 = 1071 \cdot 2^9 + 1$.
	\item $F_n$ ist nicht prim für $5 \leq n \leq 32$.
	\item Ob $F_n$ für $n = 33, 34, 35$ prim ist oder nicht, ist unbekannt. Für $n = 36$ ist wieder bekannt, dass $F_n$ nicht prim ist, ebenso für $n = 37, 38, 39$. Der nächste offene Fall ist $n = 40$.
	\item Derzeit ist von $271$ Fermatzahlen bekannt, dass sie nicht prim sind\footnote{vor 10 Jahren waren es erst $217$. Vor 40 Jahren war noch unsicher, je entscheiden zu können, ob $F_{17}$ prim ist oder nicht.}. Deren größte ist $F_{2747497}$; sie hat den Teiler $57 \cdot 2^{2747499}$. (Stand Mai 2013)
	\item Wie gesagt ist $F_{20}$ nicht prim, aber es ist kein Primteiler von $F_{20}$ bekannt; das gleiche für $F_{24}$.
\end{enumerate}

Offene (unlösbare?) Probleme:
\begin{enumerate}[1)]
	\item Gibt es unter den $F_n$ nur endlich viele Primzahlen? Man kennt bisher nur die fünf in (1).
	\item Sind unendlich viele $F_n$ nicht prim?
	\item Sind alle $F_n$ quadratfrei? Oder wenigstens unendlich viele?
\end{enumerate}

\begin{falko} \label{F7.1}
	Es gilt $F_{n+1} - 2 = \prod_{k=0}^{n} F_k$. Insbesondere sind $F_n,F_m$ für $m > n$ teilerfremd.
\end{falko}

Für welche ungeraden $n$ ist $n$-Teilung des Kreises mit Zirkel und Lineal möglich? Bisheriger Rekord:\\
$n = 3 \cdot 15 \cdot 17 \cdot 257 \cdot 65537 = F_5 - 2 = 2^{32} - 1 = 4.294.967.295$.

\minisec{Bemerkung}
	Aus F\ref{F7.1} folgt, dass es unendlich viele Primzahlen gibt: Jedes $F_n$ hat einen Primteiler $q_n$. Nach F\ref{F7.1} ist aber $q_n \neq q_m$ für $n \neq m$. \\
	Für die $n$-te Primzahl $p_n$ gilt $p_n \leq F_{n-2} = 2^{2^{n-2}}+1$ (sehr schlechte Abschätzung, aber nicht schlechter als die in §\ref{sec:para1} aus dem Beweis von Euklid: $p_n \leq 2^{2^{n-1}}$.)
	
\begin{falko}[Pépin-Test] \label{F7.2}
	Sei $n \geq 2$ und $g$ sei eine ganze, zu $F_n = 2^{2^n} + 1$ teilerfremde Zahl mit $\leg{g}{F_n}_J = -1$. Dann sind äquivalent: \begin{enumerate}[(i)]
		\item $F_n$ ist prim \index{Pépin-Test}
		\item $g^{\frac{F_n-1}{2}} \kon -1 \modu F_n$
		\item $\ord(g \modu F_n) = F_n - 1 = 2^{2^n}.$
	\end{enumerate}
\end{falko}

\minisec{Bemerkung}
	Den Test kann man zum Beispiel mit $g = 3, g= 5, g=10$ anwenden. Denn $F_n = 2^{2^n} + 1 \kon 2^{2 \cdot 2^{n-1}} \kon 1 + 1 \kon 2 \modu 3$ und $\leg{3}{F_n} = \leg{F_n}{3} = \leg{2}{3} = -1$. \\
	$F_n = 2^{4 \cdot 2^{n-2}} + 1 \kon 1 + 1 \kon 2 \modu 5$, $\leg{5}{F_n} = \leg{F_n}{5} = \leg{2}{5} = -1$, $(10, F_n) = 1$ und \\ $\leg{10}{F_n} = \leg{2}{F_n} \leg{5}{F_n} = \leg{5}{F_n} = -1$.

\minisec{Bemerkung}
	$2^k - 1$ prim $\Rightarrow k$ prim (Aufgabe 1).
	\[ M_p := 2^p - 1, p \text{ prim} \]

\begin{enumerate}[(1)]
	\item $M_2 = 3, M_3 = 7, M_5 = 31, M_7 = 127$ sind \bet{Mersennesche Primzahlen}. \index{Mersennesche Primzahl}
	\item $M_{11}$ ist keine Primzahl: $M_{11} = 2^{11} - 1 = 2047 = 23 \cdot 89$.
	\item $M_{13}, M_{17}, M_{19}$ sind prim.
	\item Für $23 \leq p \leq 100$ ist nur $M_{31}, M_{61}, M_{89}$ prim.
	\item Mersenne hat fünf Fehler gemacht: die Primzahlen $M_{61}, M_{89}, M_{107}$ übersehen und $M_{67}, M_{257}$ als prim behauptet.
	\item Für $100 \leq p \leq 257$ sind nur $M_{107}, M_{127}$ prim. \\
	$M_{127}$ bis 1951 größte bekannte Primzahl. 30. Januar 1952: $M_{521}$ ist prim! Zwei Stunden später: $M_{607}$ auch! \\
	$p$ Mersennesche Primzahl $\Rightarrow M_p$ prim? \\
	$M_3, M_7, M_{31}, M_{127}$ prim, aber $M_{8191} = M_{M_{13}}$ keine Primzahl.
	\item Für $p \leq 12000$ sind unter den ca. 1250 vielen $M_p$ genau 23 Primzahlen. $M_{11213}$ ist die größte.
	\item Bislang sind 48 Mersennesche Primzahlen bekannt (Stand Oktober 2014); die größte ist $M_{57885161}$. Sie hat 17.425.170 Stellen.
\end{enumerate}

\setcounter{countfalko}{3}
\begin{falko} \label{F7.4}
	Sei $p$ eine Primzahl mit $p \kon 3 \modu 4$. Dann gilt:
	\[2p + 1 \text{ Primzahl } \Leftrightarrow 2p+1 \text{ teilt } M_p. \]
	Ist also $2p+1$ eine Primzahl mit $p$ wie oben und $p \neq 3$, so ist $M_p$ keine Primzahl.
\end{falko}

\minisec{Bemerkung}
	Jeder Primteiler $q$ von $M_p$, $p \geq 3$, hat die Gestalt $q = 2kp + 1$.
	
\begin{falko}[Lucas-Test] \label{F7.5}
	Definiere Folge natürlicher Zahlen $s_1, s_2, \dots$ durch $s_1 = 4, s_{n+1} = s_n^2 - 2$. Dann gilt: \index{Lucas-Test}
	\[ M_p \text{ prim } \Leftrightarrow s_{p-1} \kon 0 \modu M_p \]
\end{falko}

\setcounter{countfalko}{2}
\begin{falko} \label{F7.3}
	Sei $p$ ein Primteiler einer Fermatschen Zahl $F_n$ bzw. einer Mersenne-Zahl $M_q$. Genau dann geht auch $p^2$ in $F_n$ bzw. $M_q$ auf, wenn die Bedingung \marginnote{06.01. \\ \ [21]}
	\begin{equation}
		2^{p-1} \kon 1 \modu p^2 \label{eq_W}
	\end{equation}
	erfüllt ist.
\end{falko}

Eine Primzahl, die \eqref{eq_W} erfüllt, heißt eine \Index{Wieferich-Primzahl}. Bisher sind nur die Wieferich-Primzahlen 1093 und 3511 bekannt; bis $10^{15}$ gibt es keine weiteren. Die Primzahlen $p = 1093$ und $p = 3511$ kommen aber wegen $1093 = 1 + 273 \cdot 2^2$ und $3511 = 1 + 1755 \cdot 2^1$ nicht als Primteiler einer Fermatzahl $F_n$ in Frage. Und auch nicht als Primteiler einer Mersenne-Zahl $M_q$, denn sonst wäre $q$ -- wegen $2^q \kon 1 \modu p$ -- ein Primteiler von $p-1$, also wegen $1093 - 1 = 2^2 \cdot 3 \cdot 7 \cdot 13$ und $3511 - 1 = 2 \cdot 3^3 \cdot 5 \cdot 13$ eine der Primzahlen $q = 2,3,5,7$ oder $13$. Für alle diese $q$ ist aber $M_q$ prim (und verschieden von 1093 und 3511).

\minisec{Bemerkung}
	Die \Index{Fermatsche Vermutung} besagt: Die Gleichung
	\begin{equation}
		X^n + Y^n = Z^n \label{eq_fermat}
	\end{equation}
	besitzt für $n \geq 3$ keine Lösung in natürlichen Zahlen, d.h. es gibt kein Tripel $(a,b,c) \in \NN^3$ mit $a^n+b^n=c^n$. Um dies für alle $n \geq 3$ zu etablieren, genügt es offenbar zu zeigen, dass \eqref{eq_fermat} im Falle einer Primzahl $n = p \geq 3$ keine Lösung $(a,b,c) \in \NN^3$ besitzt sowie im Falle $n = 4$. Den Fall $n = 4$ konnte bereits Fermat erledigen (F\ref{F7.3} auf Seite~\pageref{F7.3}). \\
	Im Falle $n = p \geq 3$ prim besagt der so genannte erste Fall der Fermatschen Vermutung, dass die Gleichung \eqref{eq_fermat} keine Lösung $(a,b,c) \in \NN^3$ hat, bei der keine der Zahlen $a,b,c$ durch $p$ teilbar ist. Dies, so konnte Wieferich zeigen, gilt unter der Voraussetzung, dass $p$ die Bedingung $2^{p-1} \not\kon 1 \modu p^2$ erfüllt.

\subsection{Zur Bedeutung der Mersenneschen Primzahlen}
\begin{defn}[vollkommene Zahl]
	Eine natürliche Zahl $n$ heißt $\Index{vollkommen}$, wenn sie gleich der Summe ihrer natürlichen Teiler $d < n$ ist. Wir definieren für $n \in \NN$ die Summation über alle natürlichen Teiler von $n$:
	\[ \sigma(n) := \sum_{d \mid n} d \]
\end{defn}

$n \in \NN$ ist genau dann vollkommen, wenn $\sigma(n) = 2n$. Ist $\sigma(n) < 2n$, so heißt $n$ \Index{defizient}. Andernfalls heißt $n$ \Index{abundant}.
\begin{enumerate}[(1)]
	\item $\sigma(p^k) = 1 + p + p^2 + \dots + p^k = \frac{p^{k+1}-1}{p-1} < 2p^k$, defizient
	\item $(n_1,n_2) = 1 \Rightarrow \sigma(n_1n_2) = \sigma(n_1) \sigma(n_2)$ \\
	$d_1 \mid n_1, d_2 \mid n_2 \Rightarrow d_1 d_2 \mid n_1 n_2$ \\
	$d \mid n_1 n_2 \Rightarrow d=d_1d_2$ mit $d_i \mid n_i$
	$\sum\limits_{d \mid n_1n_2} d = \sum\limits_{d_1 \mid n_1, d_2 \mid n_2} d_1d_2 = \sum\limits_{d_1 \mid n_1} d_1 \cdot \sum\limits_{d_2 \mid n_2} d_2$
\end{enumerate}

\setcounter{countfalko}{5}
\begin{falko}	\label{F7.6}
	Hat $n$ die Gestalt $n = 2^{k-1} (2^k-1)$ und ist dabei $2^k - 1$ eine Primzahl, so ist $n$ vollkommen.
\end{falko}

\begin{falko}	\label{F7.7}
	Die geraden vollkommenen Zahlen $n$ sind genau die Zahlen der Gestalt $n = 2^{p-1} M_p$ mit $M_p$ prim.
\end{falko}

\minisec{Problem (ungelöst)}
	Gibt es ungerade vollkommene Zahlen?
	
\minisec{Bemerkung}
	Man weiß: $n$ ungerade mit weniger als 8 Primteilern ist nicht vollkommen. \\
	Jedes ungerade $n \leq 10^{200}$ nicht vollkommen.
\newpage				% Elliptische PDGLs
\section{Multiplikative zahlentheoretische Funktionen}
\label{sec:para8}
	Jede Funktion $f \colon \NN \rightarrow \CC$ heiße eine \Index{zahlentheoretische Funktion}. Manchmal ist der Definitionsbereich auch $\NN_0$ oder $\ZZ$ anstelle von $\NN$.

\begin{defn}[multiplikative zahlentheoretische Funktion]	\label{def_8.1}
	$f$ wie oben heißt (zahlentheoretisch) \Index{multiplikativ}, wenn $f \neq 0$ und $f(n_1n_2) = f(n_1)f(n_2)$, falls $(n_1,n_2) = 1$. (daraus folgt stets $f(1) = 1$.)
\end{defn}

\minisec{Beispiel}
	\[ \begin{array}{c}
	\varphi(n) \text{ Eulersche } \varphi\text{-Funktion} \\ 
	\sigma(n) = \sum\limits_{d \mid n} d \\ 
	\tau(n) = \sum\limits_{d \mid n} 1 = \#\{d \in \NN : d \mid n\}
	\end{array} \]

\minisec{Bemerkungen}
	\begin{enumerate}[1)]
		\item $f$ multiplikativ $\Leftrightarrow f(n) = \prod\limits_{n} f(p^{w_p(n)})$ für alle $n \in \NN$
		\item Eine multiplikative Funktion ist festgelegt durch ihre Werte auf den Primzahlpotenzen. Ordnet man umgekehrt jeder Primzahlpotenz $> 1$ eine Zahl $\neq 0$ zu, so existiert genau eine multiplikative Funktion mit den vorgegebenen Werten.
		\item $f, g$ multiplikativ $\Rightarrow fg$ multiplikativ
		\item $f \neq 0$ heißt vollständig multiplikativ, wenn $f(n_1n_2) = f(n_1)f(n_2)$ für alle $n_1, n_2 \in \NN$. Dann gilt $f(p^k) = f(p)^k$.
	\end{enumerate}
	
\begin{falko}[Summatorische Funktion] \label{F8.1}
	Ist $f$ multiplikativ, so auch die zugehörige \Index{summatorische Funktion} $g$, definiert durch \marginnote{09.01. \\ \ [22]}
	\[ g(n) = \sum\limits_{d \mid n} f(d) \]
\end{falko}

\minisec{Beispiel (Teilersummenfunktionen)}
	\[ \sigma_\alpha(n) = \sum\limits_{d \mid n} d^\alpha = \prod\limits_{p} \frac{p^{\alpha(w_p(n)+1)}-1}{p^\alpha - 1}, \quad \alpha \in \CC \index{Teilersummenfunktion}\]

\begin{falko} \label{F8.2}
	Ist $f$ multiplikativ, so gilt für jedes $n = \prod_p p^{w_p(n)}$ die Formel
	\[ \sum\limits_{d \mid n} f(d) = \prod\limits_{p \mid n} (1+f(p) + f(p^2) + \dots + f(p^{w_p(n)})) \]
\end{falko}

\begin{defn}[Möbiusfunktion] \label{def_8.2}
	Wir definieren die multiplikative \Index{Möbiusfunktion} $\mu$ durch (vgl. Bemerkung 2 zur Definition~\ref{def_8.1}):
	\[ \mu(p^\nu) = \begin{cases}
		-1 & \text{für } \nu = 1 \\
		0 & \text{für } \nu > 1 \\
		1 & \text{für } \nu = 0 \\
	\end{cases} \]
	Also für $n \in \NN$ mit Primfaktorzerlegung $n = p_1^{k_1} \cdots p_r^{k_r}$:
	\[ \mu(n) = \begin{cases}
		(-1)^r, & \text{falls } k_1 = \dots = k_r = 1, \text{ d.h. } n \text{ quadratfrei} \\
		0, & \text{sonst}
	\end{cases} \]
\end{defn}

\begin{falko} \label{F8.3}
	Setze $\varepsilon(n) := \sum\limits_{d \mid n} \mu(d)$, d.h. $\varepsilon$ sei die zur Möbiusfunktion gehörige summatorische Funktion. Dann gilt:
	\[ \varepsilon(n) = \begin{cases}
		1 & \text{für } n=1 \\
		0 & \text{für } n \neq 1
	\end{cases} \]
	Also ist $\varepsilon$ die charakteristische Funktion zur Teilmenge $\{1\}$ von $\NN$. $\varepsilon$ ist vollständig multiplikativ (im Gegensatz zu $\mu$).
\end{falko}

\begin{falko} \label{F8.4}
	Ist $f$ multiplikativ, so gilt
	\[ \sum\limits_{d \mid n} \mu(d) f(d) = \prod\limits_{p \mid n} (1-f(p)) \]
\end{falko}

\begin{defn}[Dirichlet-Faltung] \label{def_8.3}
	Für zahlentheoretische Funktionen $f,g$ definiere die \Index{Dirichlet-Faltung} $f * g$ durch
	\[ (f * g)(n) = \sum\limits_{d \mid n} f(d) g\enbrace*{\frac{n}{d}} = \sum\limits_{\substack{a,b \in \NN \\ ab = n}} f(a) g(b) \]
	Es gelten $f * g = g*f, (f*g)*h = f*(g*h)$ sowie $\varepsilon*f = f*\varepsilon = f$.
\end{defn}

\minisec{Bemerkung}
	Die Menge $R$ der zahlentheoretischen Funktionen $f \colon \NN \rightarrow \CC$ bildet mit $+$ und $*$ einen kommutativen Ring mit Einselement $\varepsilon$.

\begin{falko} \label{F8.5}
	$f,g$ multiplikativ $\Rightarrow f * g$ multiplikativ.
\end{falko}

\begin{satz}[Möbius-Umkehrformeln] \label{satz_8.1}
	Für zahlentheoretische Funktionen $f,g$ sind folgende Aussagen äquivalent: \index{Möbius-Umkehrformeln} \begin{enumerate}[(i)]
		\item $g(n) = \sum\limits_{d \mid n} f(d)$, d.h. $g$ ist summatorische Funktion zu $f$, d.h. $g = f * i_0$
		\item $f(n) = \sum\limits_{d \mid n} \mu\enbrace*{\frac{n}{d}} g(d)$, d.h. $f = g * \mu$.
	\end{enumerate}
\end{satz}

\minisec{Korollar}
	Jede zahlentheoretische Funktion $g$ ist summatorische Funktion zu genau einer Funktion $f$. Dabei gilt:
	\[ g \text{ multiplikativ } \Leftrightarrow f \text{ multiplikativ} \]
	Ist $g$ multiplikativ so gilt für $f$ außerdem noch die Formel
	\[ f \enbrace*{\prod_p p^{\nu_p}} = \prod\limits_{p \mid n} (g(p^{\nu_p}) - g(p^{\nu_p - 1})) \]
	
\minisec{Bemerkung}
	Zu jeder zahlentheoretischen Funktion $f$ mit $f(1) \neq 0$ gibt es eine (eindeutig bestimmte) Funktion $h$ mit $f * h = h * f = \varepsilon$.
	
\minisec{Bemerkung}
	Nach dem Korollar zu Satz~\ref{satz_8.1} gibt es zu $i(n) = n$ eine eindeutig bestimmte Funktion $\varphi\colon \NN \rightarrow \CC$ mit \begin{enumerate}[(1)]
		\item $\sum\limits_{d \mid n} \varphi(d) = n$, und es gilt weiter
		\item $\varphi$ ist multiplikativ
		\item $\varphi(n) = \sum\limits_{d \mid n} \mu\enbrace*{\frac{n}{d}}d = \sum\limits_{d \mid n} \mu(d) \frac{n}{d} = n \sum\limits_{d \mid n} \frac{\mu(d)}{d}$
		\item $\varphi\enbrace*{\prod_p p^{\nu_p}} = \prod\limits_{p \mid n} \enbrace*{p^{\nu_p} - p^{\nu_p - 1}}$ bzw.1
		$\varphi(n) = \prod\limits_{p \mid n} \enbrace*{p^{w_p(n)} - p^{w_p(n)-1}} = n \prod\limits_{p \mid n} \enbrace*{1 - \frac{1}{p}}$, speziell \\
		$\varphi(p^\nu) = p^\nu - p^{\nu-1}, \nu \geq 1$
		\item $\varphi(n) = \# \{k \in \NN : 1 \leq k \leq n \text{ mit } (k,n)=1\}$
	\end{enumerate}
	
\begin{falko} \label{F8.6}
	Sei $f$ multiplikativ. Dann gilt:
	\[ f \text{ vollständig multiplikativ } \Leftrightarrow \mu f * f = \varepsilon \]
\end{falko}
\newpage				% Hölder- und Sobolevräume
\section{Der Vier-Quadrate-Satz}
\label{sec:para9}

\begin{satz}[Satz von Lagrange] \label{satz_langrange}
	Jede natürliche Zahl ist Summe von vier Quadraten, d.h. jedes $n \in \NN$ hat eine Darstellung der Gestalt \marginnote{23.01. \\ \ [26]}
	\begin{equation}
		n = x_1^2 + x_2^2 + x_3^2 + x_4^2 \text{ mit } x_i \in \ZZ \label{eq_4quad}
	\end{equation}
\end{satz}

\begin{defn}
	Für jedes $n \in \NN$ sei $r_4(n)$ die Anzahl der Quadrupel $(x_1,x_2,x_3,x_4) \in \ZZ^4$ mit \eqref{eq_4quad}.
\end{defn}

\begin{lemma}[Formel von Jacobi]
	Es gilt
	\begin{equation}
		r_4(n) = \begin{cases}
			8\sigma(n), & \text{falls } n \text{ ungerade} \\
			3\cdot 8 \sigma(u), & \text{falls } n \text{ gerade}
		\end{cases} \label{eq_jacobiformel}
	\end{equation}
	wobei $u$ den ungeraden Teil von $n$ bezeichnet, also $u = n \cdot 2^{-w_2(n)}$, und $\sigma$ die Teilersummenfunktion.
\end{lemma}

\subsection{Der Drei-Quadrate-Satz}

\begin{falko} \label{F10.1}
	Ist $n$ eine natürliche Zahl der Gestalt
	\begin{equation}
		n = 4^j \cdot (8k + 7) \text{ mit } j,k \in \NN_0, \label{eq_F10.1}
	\end{equation}
	so ist $n$ keine Summe von drei Quadraten in $\ZZ$.
\end{falko}

\minisec{Beweis}
Die Quadrate in $\ZZ/8\ZZ$ sind $0, 1$ und $4$, also ist $x^2+y^2+z^2 \not\kon 7 \modu 8$ für alle $x,y,z \in \ZZ$. Somit ist kein $n$ von der Gestalt $n = 8k+7$ eine Summe von drei Quadraten in $\ZZ$. 
Jetzt Induktion nach $j$: \\
Sei $j \geq 1$, dann ist $n = 4m$, und nach Induktionsanfang ist $m$ keine Summe von drei Quadraten. \\
Angenommen, es ist $4m = n = x^2+y^2+z^2$, dann folgt $x^2+y^2+z^2 \kon 0 \modu 4$. Dann sind $x,y$ und $z$ gerade, und es folgt $m = \leg{x}{2}^2 + \leg{y}{2}^2 + \leg{z}{2}^2 \quad \lightning$

\begin{satz}[Drei-Quadrate-Satz] \label{3quadsatz}
	Jede natürliche Zahl $n$, die nicht von der in F\ref{F10.1} genannten Gestalt \eqref{eq_F10.1} ist, ist eine Summe von drei Quadraten in $\ZZ$, d.h. $n$ hat eine Darstellung der Gestalt
	\begin{equation}
		n = x_1^2 + x_2^2 + x_3^2 \text{ mit } x_i \in \ZZ \label{eq_10_1}
	\end{equation}
	Ist zudem $n$ nicht durch $4$ teilbar, also $m \kon 1,2,3,5$ oder $6 \modu 8$, so existiert stets auch eine primitive Darstellung der Gestalt \eqref{eq_10_1}, d.h. eine mit
	\begin{equation}
		\ggT(x_1,x_2,x_3) = 1 \label{eq_10_2}
	\end{equation}
\end{satz}

Der Drei-Quadrate-Satz liegt anscheinend viel tiefer als der Vier-Quadrate-Satz, der aus dem Drei-Quadrate-Satz unmittelbar folgt: Sei $n \in \NN$ beliebig. Wir können annehmen, dass $n$ die Gestalt $n = 4^jm$ mit $m \kon 7 \modu 8$ hat (denn andernfalls ist $n$ ja schon eine Summe von drei Quadraten in $\ZZ$). Dann ist $m - 1 \kon 6 \modu 8$, also ist $m-1$ eine Summe von drei Quadraten und somit $m$ eine Summe von vier Quadraten. Aber dann ist auch $n = 4^jm$ eine Summe von vier Quadraten. \\
Der Drei-Quadrate-Satz folgt -- allerdings ohne den Zusatz \eqref{eq_10_2} -- aus einem allgemeinen Prinzip der Algebraischen Zahlentheorie ($\rightarrow$ Lokal-Global-Prinzip von Hasse). Was eine möglichst direkte Begründung des Drei-Quadrate-Satzes (ohne den Zusatz \eqref{eq_10_2}) angeht, so sind derzeit kursierende Beweise insofern "nicht elementar", als sie (anders als Gauß) den nachstehenden Satz von Dirichlet benutzen, dessen Beweis essentiell auf Methoden der komplexen Analysis beruht.

\begin{satz}[Satz von Dirichlet] \label{satz_dirichlet}
	Für beliebige $m > 1$ aus $\NN$ und jedes zu $m$ prime $a \in \ZZ$ gibt es unendlich viele Primzahlen $p$ mit $p \kon a \modu m$.
\end{satz}

\minisec{Bemerkungen}
\begin{enumerate}[1)]
	\item Auf den Zusatz \eqref{eq_10_2} im Drei-Quadrate-Satz legte schon Legendre Wert. Man beachte: Während z.B. $45 = 6^2+3^2$ im Wesentlichen die einzige Darstellung von $45$ als Summe von zwei Quadraten ist (und diese nicht primitiv ist), hat $45$ neben $45 = 6^2+3^2+0^2$ auch die primitive Darstellung $45 = 5^2+4^2+2^2$ als Summe von drei Quadraten.
	\item Gauß hat auch über die Anzahl $r_3(n)$ der Darstellungen $n = x_1^2 + x_2^2+x_3^2$ bzw. die Anzahl $R_3(n)$ der entsprechenden primitiven Darstellungen interessante Aussagen gemacht.
	\item Unabhängig von Gauß erhält man durch einfaches Abzählen der $(x_1,x_2,x_3,x_4)$ mit $n = x_1^2+x_2^2+x_3^2+x_4^2$ für $x_1 = 0, x_1=1, x_2 =2, \dots$ die banale Formel
	\begin{equation}
		r_4(n) = r_3(n) + 2r_3(n-1^2) + 2r_3(n-2^2) + 2r_3(n-3^2)+ \cdots \label{eq_10_3}
	\end{equation}
	Mit der Formel von Jacobi für $r_4(n)$ liefert \eqref{eq_10_3} eine gute Rekursionsformel für die $r_3(n)$. Aber dieser Rekursionsformel lässt sich nicht entnehmen -- jedenfalls nicht auf einfache Weise --, dass stets $r_3(n) \neq 0$, wenn $n$ nicht von der Gestalt $n = 4^j(8k+7)$ ist.
	\item In einer Extra-Vorlesung am 03. Februar 2015, 12:00 Uhr, wird Prof. Lorenz die Idee des Gaußschen Beweises für den Drei-Quadrate-Satzes vorstellen.
\end{enumerate}
\newpage				% Schwache Lösungen von elliptischen Gleichungen
% % % % % % 03. Juni
\section{Variationsrechnung und nichtlineare Gleichungen}
\label{sec:para10}
	Das Lösen einer PDGL ist oft gleichbedeutend mit der Minimierung von Energie. In der Physik sind partielle Differentialgleichungen oft eine Konsequenz aus dem Prinzip der minimalen Energie. \\
	Wir betrachten folgende Situation:
	\begin{itemize}
		\item \Index{Lagrange-Dichte} $L\colon \RR^n \times \RR \times \Omega \rightarrow \RR, (p,z,x) \mapsto L(p,z,x)$ (zur Vereinfachung als glatt vorausgesetzt, Ableitungen: $L_p,L_z,L_x$)
		\item Energie $E[u] = \int_\Omega L(\nabla u(x),u(x),x)dx$, definiert für (schwach) differenzierbare Funktionen $u\colon \Omega \rightarrow \RR$.
		\item \Index{Gâteaux-Differential} von $E$ in Richtung $v\colon \Omega \rightarrow \RR$:
		\[ \der_u E[u](v) = \frac{d}{dt} E[u+tv] = \int_{\Omega} L_p(\nabla u,u(x),x) \cdot \nabla v + L_z(\nabla u, u, x) v dx \]
	\end{itemize}
	Betrachte das Minimierungsproblem
	\[ \min_{u \colon \Omega \rightarrow \RR} E[u] \text{ in Abhängigkeit von } u = g \text{ auf } \der \Omega \]
	Wir nehmen an, es besitzt einen glatten Minimierer $u^*$. Dann ist 0 ein Minimierer von $t \mapsto E[u^*+tv]$ für jedes $v \in C_c^\infty(\Omega)$, d.h. $\der_u E[u^*](v) = 0$. Partielle Integration liefert
	\[ 0 = \int_{\Omega} v(-\diver L_p(\nabla u^*,u^*,x) + L_z(\nabla u^*,u^*,x)) dx \]
	für alle $v \in C_c^\infty (\Omega)$, d.h. $u^*$ löst die PDGL
	\[ 0 = -\div L_p(\nabla u(x),u(x),x) + L_z(\nabla u(x),u(x),x) \text{ auf } \Omega \quad \text{ mit } u = g \text{ auf } \der \Omega \]
	
\begin{bsp} \label{bsp_75}
\marginnote{[75]} \begin{tabular}{cccc}
Langrange-Dichte $L(p,z,x)$ & Energie $E[u]$ & PDGL &  \\ 
$\frac{1}{2} |p|^2$ & $\int_\Omega \frac{1}{2} |\nabla u|^2 dx$ & $\Delta u = 0$ & (Laplace-Gleichung) \\ 
$\frac{1}{2} p^T A(x)p - zf(x)$ & $\int_\Omega \frac{1}{2} \nabla u^T A \nabla u - fu dx$ & $-\diver(A\nabla u) = f$ & (allgemeine Poisson-Gleichung) \\ 
$\frac{1}{2} |p|^2 - F(z)$ & $\int_\Omega \frac{1}{2} |\nabla u|^2 - F(u) dx$ & $-\Delta u = F'(u)$ & (nichtlineare Poisson-Gleichung) \\ 
$\sqrt{1+|p|^2}$ & $\int_\Omega \sqrt{1+ |\nabla u|^2} dx$ & $\diver \frac{\nabla u}{\sqrt{1+ |\nabla u|^2}} = 0$ & Minimalflächengleichung \\ 
\end{tabular}
\end{bsp}

\mbox{} \\
Wir führen nun ein paar Hilfsmittel aus der Funktionalanalysis ein und beweisen dann die Existenz von Minimierern (und damit von Lösungen der Differentialgleichungen) für eine umfassende Klasse nichtlinearer Energien.

\begin{defn}[Dualraum] \label{def_76}
	Der \Index{Dualraum} $X'$ \marginnote{[76]} eines Banachraums $X$ ist der Raum der beschränkten linearen Funktionale $f\colon X \rightarrow \RR$ auf $X$.
\end{defn}

\begin{defn}[schwache Konvergenz, schwach-*-Konvergenz] \label{def_77}
	Eine Folge $x_k \in X$ \bet{konvergiert schwach} gegen $x \in X$, falls $f(x_k) \rightarrow f(x)$ für alle $f \in X'$. Schreibe: $x_k \weak x$. \marginnote{[77]} \index{schwache Konvergenz}\\
	Eine Folge $f_k \in X'$ \bet{konvergiert schwach-*} gegen $f \in X'$, falls $f_k(x) \rightarrow f(x)$ für alle $x \in X$. Schreibe: $f_k \weaks f$. \index{schwach-*-Konvergenz} \\
	Offensichtlich folgt aus der Konvergenz in $X$ bzw. $X'$ auch die schwache Konvergenz bzw. die schwach-*-Konvergenz.
\end{defn}
	
\begin{thm}[Schwach-*-Kompaktheit] \label{thm_78}
	Die Einheitskugel (und damit jede abgeschlossene Teilmenge)\marginnote{[78]} eines separablen Banachraums $X$ ist schwach-*-folgenkompakt, d.h. jede Folge enthält eine konvergente Teilfolge.
\end{thm}

\minisec{Beweis}
	siehe z.B. Alt, \glqq Lineare Funktionalanalysis\grqq, S. 229.
	
\begin{defn}[Reflexivität] \label{def_79}
	Ein Banachrau $X$ heißt \Index{reflexiv}, falls er isometrisch isomorph ist zu seinem Bidualraum $(X')'$. \marginnote{[79]} \\
	Hinweis: Nach dem vorigen Satz besitzen beschränkte Folgen in $X$ eine schwach konvergente Teilfolge, falls $X$ separabel ist.
\end{defn}

\begin{thm}[Sobolevräume sind reflexiv] \label{thm_80}
	Sei $\Omega$ offen, beschränkt und mit Lipschitzrand sowie $k \in \NN_0$. Für $p \in (1,\infty)$ ist $W^{k,p}(\Omega)$ separabel und reflexiv. \marginnote{[80]}
\end{thm}
	
\minisec{Beweis}
	siehe z.B. Alt, \glqq Lineare Funktionalanalysis\grqq, S. 234.
	
% % % % % 6. Juni
\begin{bem} \label{bem_81}
	Mit der Reflexivität von $L^p(\Omega)$ für $p \in (1,\infty)$ \marginnote{6. Jun. \\ \ [81]} und der Hölder-Ungleichung ist leicht zu sehen, dass $L^p(\Omega) \subseteq (L^{p^*}(\Omega))'$. Gleichermaßen erkennt man $L^{p^*}(\Omega) \subseteq (L^p(\Omega))'$. Tatsächlich folgt mit dem Satz von Radon-Nikodým:
	\[ (L^p(\Omega))' = L^{p^*}(\Omega) \]
\end{bem}
	
\begin{thm} \label{thm_82}
	Sei \marginnote{[82]} \begin{itemize}
	\item $L(p,z,x)$ konvex in $p$
	\item $L(p,z,x)$ halbstetig von unten in $z$
	\item $L(p,z,x) \geq \alpha |p|^q - \beta$ für gewisse $\alpha, \beta > 0, q \in (1,\infty)$.
	\end{itemize}
	Falls $g \in W^{1,q}(\Omega)$, dann hat $E$ einen Minimierer in $\{u \in W^{1,q}(\Omega) : u = g \text{ auf } \der \Omega \}$.
\end{thm}
	
\minisec{Beweis}
	\glqq Direkte Methode der Variationsrechnung\grqq \begin{enumerate}[1.]
		\item Weder gilt $E \equiv \infty$, noch ist $E$ nach unten unbeschränkt.
		\item Betrachte eine minimierende Folge $u_k$ mit $E[u_k] \rightarrow \inf_u E[u]$ monoton.
		\item Zeige die Kompaktheit der Folge, d.h. die Existenz einer (in einem gewissen Sinn) konvergenten Teilfolge $u_k \rightarrow u^*$ für ein $u^*$. \\
		Es ist $E[u] \geq \alpha ||\nabla u||_{L^q(\Omega)}^q - \beta |\Omega|$, daher existiert ein $C > 0$ mit $||\nabla u_k||_{L^q(\Omega)} \leq C$ für alle $k$. Mit der Poincaré-Ungleichung folgt $||u_k||_{W^{1,q}(\Omega)} \leq C$ für eine (andere) Konstante $C \geq 0$. Da $W^{1,q}(\Omega)$ reflexiv ist, konvergiert eine Teilfolge $u_k$ schwach gegen ein $u^* \in W^{1,q}(\Omega)$.
		\item Zeige mit $u_k$, dass $E$ halbstetig von unten ist, d.h. $E[u] \leq \liminf_{k \rightarrow \infty} E[u_k]$. \\
		Aufgrund der kompakten Einbettung $W^{1,q}(\Omega) \hookrightarrow L^q(\Omega)$ gilt $u_k \rightarrow u^*$ im starken Sinne in $L^q(\Omega)$ für eine Teilfolge und demach sogar punktweise fast überall nach dem Entfernen einer anderen Teilfolge. Mit dem Satz von Egoroff finden wir sogar für jedes $\varepsilon > 0$ eine messbare Menge $\Omega_\varepsilon \subseteq \Omega$ mit $|\Omega \setminus \Omega_\varepsilon| < \varepsilon$ und $u_k \rightarrow u^*$ gleichmäßig in $\Omega_\varepsilon$. \\
		O.B.d.A. sei $L \geq 0$. Da $L$ konvex in $p$, folgt
		\[ \liminf_{k \rightarrow \infty} \int_{\Omega} L(\nabla u_k,u_k,x) dx \geq \liminf_{k \rightarrow \infty} \enbrace*{ \int_{\Omega} L(\nabla u^*,u_k,x) dx + \int_{\Omega} L_p(\nabla u^*,u_k,x) \cdot (\nabla u_k - \nabla u^*) dx }. \]
		Damit haben wir
		\begin{equation}
		\begin{aligned}
			\liminf_{k \rightarrow \infty} \int_{\Omega} L(\nabla u^*,u_k,x) dx \geq \ &\liminf_{k \rightarrow \infty} \int_{\Omega} \inf_{j \geq k} L(\nabla u^*,u_j,x) + \beta dx - \beta |\Omega| \\ \notag
			\overset{(A)}{=} \ &\int_{\Omega} \liminf_{k \rightarrow \infty} L(\nabla u^*,u_k,x)+ \beta dx - \beta |\Omega| \overset{(B)}{\geq} \int_{\Omega} L(\nabla u^*,u^*,x)dx = E[u^*]
		\end{aligned}
		\end{equation}
		wobei wir bei (A) den Satz von der monotonen Kovergenz (Beppo Levi) bzw. das Lemma von Fatou und bei (B) die Halbstetigkeit von unten von $L$ benutzen. Weiter folgt:
		\begin{equation}
		\begin{aligned}
			&\int_\Omega L_p(\nabla u^*,u_k,x) \cdot (\nabla u_k,-\nabla u^*) dx \\ \notag
			\geq \ &\int_{\Omega_\varepsilon} L_p(\nabla u^*,u^*,x) \cdot (\nabla u_k,\nabla u^*) dx + \int_{\Omega_\varepsilon} (L_p(\nabla u^*,u_k,x) - L_p(\nabla u^*,u^*,x)) \cdot (\nabla u_k - \nabla u^*) dx
		\end{aligned}
		\end{equation}
		Der erste Term konvergiert wegen der schwachen Konvergenz von $u_k$ gegen 0 und der zweite wegen der Hölder-Ungleichung und der gleichmäßigen Konvergenz. Mit $\varepsilon \rightarrow 0$ folgt die Behauptung. \qed
	\end{enumerate}
\newpage			% Variationsrechnung und nichtlineare Gleichungen
% % % % % % 17. Juni
\section{Parabolische PDGL}
\label{sec:para11}
	
\begin{defn}[Wärmeleitungsgleichung]
	Die \Index{Wärmeleitungsgleichung} \marginnote{17. Jun}
	\begin{equation}
		u_t - \Delta u = 0 \label{eq_32}
	\end{equation}
	beschreibt die zeitliche Entwicklung einer diffundierenden Größe (z.B. Wärme).
	\begin{tabbing}
	\hspace{3cm}\=\kill
	Temperatur: \> $u \colon (0,T) \times \Omega \rightarrow \RR$ (in einem Stück Material $\Omega$ zur Zeit $(0,T)$)\\ 
	Leitfähigkeit: \> $a > 0$ (Materialparameter)\\ 
	Wärmekapazität: \> $\kappa > 0$ (Materialparameter)\\ 
	Wärmefluss: \> $F = -a \nabla u$ (in Richtung des negativen Temperaturgradienten)\\ 
	Nettofluss: \> $- \int_{\der V} F \cdot \nu dx$ nach $V \subset \Omega$\\ 
	Wärmeänderung: \> $\frac{d}{dt} \int_V \kappa u dx$ in $V$ 
	\end{tabbing} 
	Wärmeänderung = Nettofluss $\Rightarrow$
	\[ 0 = \frac{d}{dt} \int_V \kappa u dx + \int_{\der V} F \cdot \nu dx = \int_V \kappa u_t + \diver F dx = \int_V \kappa u_t - a \Delta u dx \]
	und \eqref{eq_32} folgt für $\kappa = a$, da $V$ beliebig ist.
\end{defn}
	
	\mbox{} \\
	Parabolische Gleichungen sind eng mit elliptischen verwandt, daher folgen wir etwa den elliptischen Methoden. Diesmal beginnen wir jedoch mit der Fundamentallösung, da Mittelwertformeln und Maximumsprinzipien dann einfacher folgen.
	
	\mbox{} \\
	Suche nach einer einfachen, radialsymmetrischen Lösung, z.B. der Form
	\[ u(t,x) = t^{-\alpha} v(t^{-\beta} r) \]
	mit $r = |x|$. Mit $\Delta = \der_r^2 + \frac{n-1}{r} \der_r$ und $y = t^{-\beta} r$ ergibt sich
	\[ -\alpha t^{-\alpha - 1} v(y) - \beta t^{-\alpha-1} y v'(y) - t^{-\alpha - 2\beta} v''(y) + \frac{n-1}{y} t^{-\alpha - 2\beta} v'(y) = 0\]
	Wähle $\beta = \frac{1}{2}$. Mit $\alpha = \frac{n}{2}$ wird dies zu
	\[ 0 = [y^{n-1} v' + \frac{1}{2} y^n v]' = [vy^{n-1} (\log v + \frac{y^2}{4})']' \]
	d.h. eine Lösung ist $v = c \exp \enbrace*{-\frac{y^2}{4}}$.
	
\begin{defn}[Fundamentallösung] \label{def_83}
	Die Funktion \marginnote{[83]}
	\[ \Phi(t,x) = \begin{cases}
		\frac{1}{(4\pi t)^{n/2}} \exp \enbrace*{ - \frac{|x|^2}{4t}} & (t > 0) \\
		0 & (t < 0) \end{cases} \]
	löst \eqref{eq_32} auf $(\RR \setminus \setnull) \times \RR^n$ und heißt \Index{Fundamentallösung} der Wärmeleitungsgleichung.
\end{defn}
	
\begin{bem} \label{bem_84}
	Zur Zeit $t > 0$ \marginnote{[84]} ist die Fundamentallösung offfensichtlich eine $n$-dimensionale Normalverteilung/Gaußsche Glockenkurve mit Mittelwert 0 und Standardabweichung $\sqrt{2t}$.
\end{bem}
	
\minisec{Beweis}
	\begin{description}
		\item[$n=1$:]
		\[ \enbrace*{ \int_R \frac{1}{\sqrt{4 \pi t}} e^{-\frac{x^2}{4t}} dx}^2 = \int_{\RR} \frac{1}{4\pi t} e^{-\frac{x^2+y^2}{4t}} dxdy \overset{PK}{=} \int_{r=0}^{\infty} 2\pi r \frac{1}{4\pi t} e^{-\frac{r^2}{4t}} dr = \left \lbrack -e^{- \frac{r^2}{4t}} \right\rbrack_{r=0}^\infty = 1 \]
		\item[$n > 1$:]
		\[ \int_{\RR^n} \frac{1}{\sqrt{4\pi t}^n} e^{- \frac{|x|^2}{4t}} dx = \prod\limits_{i=1}^{n} \int_{\RR} \frac{1}{\sqrt{4\pi t}} e^{-\frac{|x_i|^2}{4t}} dx_i = 1 \]
	\end{description}
	
\mbox{} \\
Definiere: \begin{itemize}
	\item $\Omega_T = \Omega \times (0,T]$ \\
	\item $\Gamma_T = \overline{\Omega}_T \setminus \Omega_T = \setnull \times \Omega \cup [0,T] \times \der \Omega$
	\item \Index{Wärmeball} für $t \in \RR, x \in \RR^n$ (ein Levelset der Fundamentallösung)
	\[ E(t,x,r) = \penbrace*{(s,y) \in \RR^{n+1} : s \leq t, \Phi(t-s,x-y) \geq \frac{1}{r^n}} \]
\end{itemize}

\begin{thm}[Mittelwertformel] \label{thm_85}
	Sei $u \in C^1((0,T]),C^2(\Omega,\RR)$ eine Lösung von \eqref{eq_32} in $\Omega_T$, dann gilt \marginnote{[85]}
	\begin{equation}
		u(t,x) = \frac{1}{4r^n} \int_{E(t,x,r)} u(s,y) \frac{|x-y|^2}{(t-s)^2} dyds =: U(r) \label{eq_33}
	\end{equation}
	für jedes $E(x,r,t) \subset \Omega_T$. (Beachte: Formel integriert nur über Vergangenheit!)
\end{thm}
	
\minisec{Beweis}
	O.B.d.A sei $t=0, x=0$. $U(r) = \frac{1}{4} \int_{E(0,0,r)} u(r^2 \tilde{s},r \tilde{y}) \frac{|0- \tilde{y|^2}}{(0-\tilde{s})^2} d\tilde{y} d\tilde{s}$. \\
	\begin{equation}
	\begin{aligned}
	U'(r) &= \frac{1}{4} \int_{E(0,0,1)} 2rs u_t(r^2 s,ry) \frac{y^2}{s^2} + yu_x(r^2s,ry) \frac{y^2}{s^2} dyds \\ \notag 
	&= \frac{1}{4r^{n+1}} \int_{E(0,0,r)} \nabla u(s,y) \cdot y \frac{|y|^2}{s^2} + 2u_s \frac{|y|^2}{s} dyds =: A+B.
	\end{aligned}
	\end{equation}
	Sei $\psi := \log(r^n \Phi(-s,y)) = -\frac{n}{2} \log (-urs) + \frac{|y|^2}{4^s} + n \log r$. Wegen $\Phi(-s,y) = r^{-n}$ auf $\der E(0,0,r)$ ist dort $\psi = 0$.
	\begin{equation}
	\begin{aligned}
		B &= \frac{1}{4r^{n+1}} \int_{E(0,0,r)} 4u_s y \cdot \nabla \psi dy ds \overset{P.I.}{=} -\frac{1}{4r^{n+1}} \int_{E(0,0,r)} \psi n u_s + \psi y \cdot \nabla u_s dy ds \\ \notag
		&\overset{P.I.}{=} \frac{1}{r^{n+1}} \int_{E(0,0,r)} -\psi n u_s + \underbrace{\psi_s}_{-\frac{n}{2s} - \frac{|y|^2}{4s^2}} y \cdot \nabla u dy dy \\
		&= \frac{1}{r^{n+1}} \int_{E(0,0,r)} -\psi n \Delta u - \frac{n}{2s} y \cdot \nabla u dy ds - A = \frac{1}{r^{n+1}} \int_{E(0,0,1)} n \nabla \psi \cdot \nabla u - \frac{n}{2s} y \cdot \nabla u dy ds -A = -A
	\end{aligned}
	\end{equation}
	$\Rightarrow U'(r) = 0 \Rightarrow U(r) = \lim\limits_{\rho \rightarrow 0} U(\rho) = u(0,0) \lim\limits_{\rho \rightarrow 0} \underbrace{\frac{1}{4\rho^n} \int_{E(0,0,\rho)} \frac{|y|^2}{s^2} dyds}_{=\frac{1}{4} \int_{E(0,0,1)} \frac{|y|^2}{s^2} dyds = 1}$ \qed

\begin{thm}[Starkes Maximumprinzip] \label{thm_86}
	$u \in C^1((0,T];C^2(\Omega,\RR)) \cap C(\Omega)$ sei eine Lösung von \eqref{eq_32} in $\Omega_T$, dann gilt: \marginnote{[86]}
	\[ \max_{\overline{\Omega_T}} = \max_{\Gamma_T} u \]
	Ist $u(t,x) = \max_{\overline{\Omega_T}} u$ für $(t,x) \in \Omega_T$, so ist $u$ konstant auf $\Omega_t$, falls $\Omega$ zusammenhängend ist. \\
	(Analog für Minima durch Ersetzen von $u$ mit $-u$.)
\end{thm}

\minisec{Beweis} 
	Sei $M = u(x,t)$ \marginnote{20. Jun} ein Maximum. Nach Theorem \ref{thm_85} ist $M = \frac{1}{4r^n} \int_{E(t,x,r)} u(s,y) \frac{|x-y|^2}{|t-s|^2} dyds \leq M$, da \linebreak $\int_{E(t,x,r)} \frac{|x-y|^2}{(t-s)^2} dyds = 1$. Gleichheit gilt genau dann, wenn $u \equiv M$ auf $E(t,x,r)$. Um jeden Punkt $(s,y)$ von $E(t,x,r)$ kann wieder ein Wärmeball gezeichnet werden, auf dem $u \equiv M$ ist. So kann $\Omega_t$ mit Wärmebällen überdeckt werden, auf denen $u$ jeweils gleich $M$ ist. \qed

\begin{bem}[Unendliche Ausbreitungsgeschwindigkeit] \label{bem_87}
	Sei $u \in C^1((0,T];C^2(\Omega,\RR)) \cap C(\Omega_T)$ eine Lösung des Anfangs-Randwertproblems \marginnote{[87]}
	\[ \begin{cases}
		u_t - \Delta u = 0 & \text{ auf } \Omega_T \\
		u = 0 & \text{ auf } [0,T] \times \der \Omega \\
		u = g \text{ auf } \setnull \times \Omega \end{cases} \]
	für $g \geq 0$. Aus dem Maximumsprinzip folgt, dass $u > 0$ auf $\Omega_T$, falls ein $x \in \Omega$ existiert mit $g(x) > 0$, d.h. die Anfagswert-Information ist unendlich schnell überall hingeflossen.
\end{bem}
	
\begin{defn}[Inhomogene Wärmeleitung]
	Die inhomogene Wärmeleitungsgleichung
	\begin{equation}
		u_t(t,x) - \Delta u(t,x) = f(t,x) \label{eq_34}
	\end{equation}
	hat einen Quellterm $f$. Adäquate Randbedingungen für parabolische Gleichungen sind eine Anfangsbedingung
	\begin{equation}
		u=g \text{ für } t = 0 \label{eq_35}
	\end{equation}
	und eine Dirichlet- oder Neumann-Randbedingung
	\begin{equation}
		u = h \text{ auf } \der \Omega \label{eq_36}
	\end{equation}
	oder
	\begin{equation}
		\der u / \der v = h \text{ auf } \der \Omega \label{eq_37}
	\end{equation}
\end{defn}

\begin{thm}[Eindeutigkeit] \label{thm_88}
	Sei $g \in C(\Omega), h \in C([0,T] \times \der \Omega)$, \marginnote{[88]} dann ist eine Lösung $u \in C^1((0,T];C^2(\Omega,\RR)) \cap C(\Omega_T)$ von \eqref{eq_34}, \eqref{eq_35} und \eqref{eq_36} eindeutig.
\end{thm}

\minisec{Beweis}
	Das folgt aus dem starken Maximumsprinzip für $u - \tilde{u}$, wenn $u,\tilde{u}$ zwei Lösungen sind. \qed
	
\begin{thm}[Fundamentallösung] \label{thm_89}
	Es gilt \marginnote{[89]}
	\[ \begin{cases}
		\mathcal{L}[\Phi] := \Phi_t - \Delta \Phi = 0 & \text{ auf } (0,\infty) \times \RR^n \\
		\Phi(0,x) = \delta(x) \end{cases} \]
	in dem Sinn, dass für jedes $f \in C_c^1(\RR^n)$ gilt:
	\[ \lim_{t\rightarrow 0} \int_{\RR^n} \Phi(t,x) f(x) dx = f(0) \]
\end{thm}
		
\minisec{Beweis}
	Nebenrechnung:
	\begin{equation}
	\begin{aligned}
		& \ \enbrace*{ \overbrace{\int_{\RR \setminus [-\sqrt[4]{t},\sqrt[4]{t}]} \frac{1}{\sqrt{4 \pi t}} e^{- \frac{x^2}{4t}} dx}^{=:A_t}}^2 = \int_{\RR^2 \setminus [-\sqrt[4]{t},\sqrt[4]{t}] \times [-\sqrt[4]{t},\sqrt[4]{t}]} \frac{1}{\sqrt{4 \pi t}} e^{- \frac{x^2+y^2}{4t}} dxdy \\ \notag
		\leq \ & \ \int_{r=2 \sqrt[4]{t}}^{\infty} 2\pi r \frac{1}{4\pi t} e^{- \frac{r^2}{4t}} dr = \left\lbrack -e^{-\frac{r^2}{4t}} \right\rbrack_{r = 2\sqrt[4]{t}}^\infty = e^{-\frac{1}{\sqrt{t}}}
	\end{aligned}
	\end{equation}
	\begin{equation}
		\Rightarrow \int_{\RR^n} \Phi(t,x) f(x) dx = \int_{\RR^n \setminus [-\sqrt[4]{t},\sqrt[4]{t}]^n} \Phi(t,x) f(x) dx + \int_{[-\sqrt[4]{t},\sqrt[4]{t}]^n} \Phi(t,x) f(x) dx \notag
	\end{equation}
	Der erste Summand ist betragsäßig kleiner als $A_t^n \max_\RR^n |f| \xlongrightarrow[t \rightarrow 0]{} 0$. Der zweite Summen liegt zwischen zwischen dem Minimum und dem Maximum von $f(1-A_t)$ auf $[-\sqrt[4]{t},\sqrt[4]{t}]$; beides strebt gegen $f(0)$ für $t \rightarrow 0$. \qed

% % % % 24. Juni
\begin{bem} \label{bem_90}
	Alternativ (und vielleicht näher zu unserer Vorgehensweise für elliptische Gleichungen) kann $\Phi$ \marginnote{24. Jun \\ \ [90]} aufgefasst werden als Lösung zu
	\[ \begin{cases}
		\Phi_t - \Delta \Phi = \delta & \text{ auf } \RR \times \RR^n \\
		\Phi(0,x) = 0 \end{cases} \]
\end{bem}

\mbox{} \\
Im selben Sinn wie zuvor, betrachte nun für festes $(s,y)$ die Lösung des zu \eqref{eq_34}-\eqref{eq_36} \textit{adjungierten Problems} (eine Diffusion rückwärts in der Zeit)
\begin{equation}
	\begin{cases}
		\mathcal{L}^*[G^{(s,y)}] := -G_t^{(s,y)} - \Delta G^{(s,y)} = 0 & \text{auf } (0,s) \times \Omega \\
		G^{(s,y)}(t,x) = 0 & \text{auf } \der \Omega \\
		G^{(s,y)}(t,x) = \delta(x-y) & \text{zu } t=s \end{cases} \label{eq_38}
\end{equation}
Motivation: wenn wir $G^{(s,y)}$ für alle $(s,y)$ finden, gilt (informell; formaler Beweis wie bei elliptischen Differentialgleichungen)
\begin{equation}
\begin{aligned}
	\int_{\Omega_s} (G^{(s,y)} \mathcal{L}[u] - u \mathcal{L}^*[G^{(s,y)}]) dxdt
	&= \int_{\Omega_s} \der_t (G^{(s,y)} u) + \nabla(uG_x^{(s,y)} - G^{(s,y)} u_x) dx dt \\ \notag
	&= \int_{\Omega_s} \begin{pmatrix} \der_t \\ \nabla \end{pmatrix} \cdot \begin{pmatrix} G^{(s,y)} u \\ uG_x^{(s,y)} - G^{(s,y)} u_x \end{pmatrix} dxdt \\
	&= \int_{[0,s]\times \der\Omega} u \der_v G^{(s,y)} - G^{(s,y)} \der_v u dx dt + \int_\Omega G^{(s,y)} u \bigg|_{t=s} - G^{(s,y)} u \bigg|_{t=0} dx
\end{aligned}
\end{equation}
und somit
\[ u(s,y) = \int_{\Omega_s} G^{(s,y)} f dx dt + \int_{\setnull \times \Omega} G^{(s,y)} gdx - \int_{[0,s]\times \der \Omega} \der_v G^{(s,y)} dx dt,\]
eine \textit{Greensche Darstellung} der Lösung.

\begin{bsp}[Greensche Funktion für Halbraum] \label{bsp_91}
	$G(s,y)$ für $\Omega = \{x \in \RR^n : x_n > 0\}$ kann wieder mit der Spiegelungsmethode (vgl. \ref{bsp_49}) gefunden werden:
	\[ G^{(s,y)}(t,x) = \Phi(s-t,x-y) - \Phi(s-t,x+y) \]
	erfüllt \eqref{eq_38}.
\end{bsp}
\newpage			% Parabolische PDGL
\section{Schwache Lösungen von parabolischen Gleichungen}
\label{sec:para12}
	
Erinnerung: \begin{itemize}
	\item $H_0^1(\Omega) = \{ f \colon \Omega \rightarrow \RR : \int_{\Omega} |f|^2 dx < \infty, \text{ schwache Ableitung } Df \text{ existiert}, \int_{\Omega} |Df|^2 dx, f \big|_{\der \Omega} = 0\}$
	\item Banachraum: normierter, vollständiger Vektorraum
	\item Dualraum $X^*$ zu Banachraum $X$: $X^* = \{f \colon X \rightarrow \RR : f \text{ ist linear}, ||f||_{X^1} := \sup\limits_{\substack{x \in X \\ ||x||_X \leq 1}} f(x) < \infty\}$
	\end{itemize}
	
\begin{defn}
	$H^{-1}(\Omega)$ bezeichne den Dualraum zu $H_0^1(\Omega)$. Für $f \in H^{-1}(\Omega)$ schreiben wir auch $\langle f,v \rangle$ statt $f(v)$.
	\[ ||f||_{H^{-1}(\Omega)} = \sup\limits_{\substack{u \in H_0^1(\Omega) \\ ||u||_{H_0^1(\Omega)} \leq 1}} \langle f,u \rangle \]
\end{defn}
	
\begin{thm}
	\begin{enumerate}[(i)]
		\item Sei $f \in H^{-1}(\Omega)$, dann gibt es Funktionen $f^0,f^1,\dots,f^n \in L^2(\Omega)$ mit
		\begin{equation}
			\langle f,v \rangle = \int_{\Omega} f^0 v dx + \sum\limits_{i=1}^{n} \int_{\Omega} f^i v_{x_i} dx \quad \forall v \in H_0^1(\Omega) \label{eq_stern}
		\end{equation}
		\item $||f||_{H^{-1}(\Omega)} = \inf\limits_{f \text{ erf. \eqref{eq_stern}}} \sqrt{\int_{\Omega} \sum\limits_{i=0}^{n} |f^i|^2}$
		\item $(u,v)_{L^2(\Omega)} = \int_{\Omega} uvdx = \langle u,v \rangle$ für alle $v \in H_0^1(\Omega), u \in H^{-1}(\Omega) \cap L^2(\Omega)$.
	\end{enumerate}
\end{thm}
	
\minisec{Beweis}
	\begin{enumerate}[(i)]
	\item $H_0^1(\Omega)$ ist mit Skalarprodukt $(u,v)_{H_0^1(\Omega)} = \int_{\Omega} uv + Du Dv dx$ ein Hilbertraum. Nach \ref{thm_68} kann jedes $f \in H^{-1}(\Omega)$ dargestellt werden durch ein $u_f \in H_0^1(\Omega)$, d.h. $\langle f,v \rangle = (u_f,v)_{H_0^1(\Omega)}$ für alle $v \in H_0^1(\Omega)$. Dann folgt (i) mit $f^0 = u_f$, $f^i = (u_f) _{x_i}$.
	\item Sei $f \in H^{-1}(\Omega)$ mit $\langle f,v \rangle = \int_{\Omega} g^0 v + \sum_{i=1}^n g^i v_{x_i} dx$ für $g^0,\dots,g^n \in L^2(\Omega)$. Es gilt:
	\begin{equation}
	\begin{aligned}
		\int_{\Omega} g^0 v + \sum\limits_{i=1}^{n} g^i v_{x_i} dx \quad &\hspace{-2mm}\overset{\text{Hölder}}{\leq} ||g^0||_{L^2} ||v||_{L^2} + \sum\limits_{i=1}^{n} ||g^i||_{L^2} \\
		&\hspace{-9mm}\overset{a \cdot b \leq |a||b| \text{ in } \RR^{n+1}}{\leq} \sqrt{\sum\limits_{i=0}^{n} ||g^i||_{L^2}^2} \sqrt{||v||_{L^2}^2 + \sum\limits_{i=1}^n ||v_{x_i}||_{L^2}^2} \\
		&\leq \sqrt{\int_{\Omega} \sum\limits_{i=0}^n |g^i|^2 dx} \cdot ||v||_{H_0^1(\Omega)} \label{eq_stern2}
	\end{aligned}
	\end{equation}
	Nun:
	\[ ||f||_{H^{-1}(\Omega)} = \sup\limits_{||u||_{H_0^1(\Omega)} \leq 1} \langle f,u \rangle = \sup\limits_{||u||_{H_0^1(\Omega)} \leq 1} \int_{\Omega} g^0 u + \sum_{i=1}^n g^i u_{x_i} dx \overset{\text{\eqref{eq_stern2}}}{\leq} \sqrt{\int_{\Omega} \sum\limits_{i=0}^{n} |g^i|^2 dx} \]
	und
	\[ ||f||_{H^{-1}(\Omega)} \geq \langle f,\frac{u_f}{||u_f||_{H_0^1}} \rangle = \frac{\int_{\Omega} f^0 u + \sum\limits_{i=1}^{n} f^i u_{x_i} dx}{\sqrt{\int_\Omega u_f^2 + |Du_f|^2 dx}} = \frac{\int_{\Omega} \sum\limits_{i=0}^{n} |f^i|^2 dx}{\sqrt{\int_\Omega \sum\limits_{i=0}^{n} |f^i|^2 dx}} = \sqrt{\int_{\Omega} \sum\limits_{i=0}^{n} |f^i|^2 dx} \]
	\item folgt aus (i). \qed
	\end{enumerate}
	
\begin{defn}[$L^p$-Raum]
\begin{itemize}
	\item Für einen Banachraum $X$ sei
	\[L^p((0,T);X) := \{ u \colon (0,T) \rightarrow X : u \text{ ist stark messbar und } ||u||_{L^p((0,T);X)} < \infty\} \]
	mit $||u||_{L^p((0,T);X)} = \enbrace*{ \int_{0}^{T} ||u(t)||^p dt}^{\frac{1}{p}}$ für $1 \leq p < \infty$ und $||u||_{L^\infty((0,T);X)} = \esssup_{t \in (0,T)} ||u(t)||$. \\
	(Dabei heißt $u$ stark messbar, falls eine fast überall gegen $u$ konvergente Folge von Treppenfunktionen existiert.)
	\item $C([0,T];X) := \penbrace*{ u\colon [0,T] \rightarrow X : u \text{ stetig}, ||u||_{C([0,T];X)} := \max\limits_{t \in [0,T]} ||u(t)|| < \infty}$
\end{itemize}
\end{defn}

\begin{defn}[schwache Ableitung]
	Sei $u \in L^1((0,T);X)$. $v \in L^1((0,T);X)$ heißt \Index{schwache Ableitung} von $u$ und wir schreiben $u' = v$, wenn
	\[ \int_{0}^{T} \phi'(t) u(t) dt = -\int_{0}^{T} \phi(t) v(t) dt \]
	für alle $\phi \in C_c^\infty ((0,T);\RR)$.
\end{defn}
	
\begin{bem}
	$u \in L^2((0,T);H_0^1(\Omega))$ und $u' \in L^2((0,T);H^{-1}(\Omega))$ impliziert: \begin{itemize}
		\item $u \in C([0,T];L^2(\Omega))$
		\item $||u(t)||_{L^2(\Omega)}^2 = ||u(0)||_{L^2(\Omega)}^2 + \int_0^t 2 \langle u'(t),u(t) \rangle dt$ \\
		(bzw. $\frac{d}{dt} ||u(t)||_{L^2}^2 = \frac{d}{dt} \int_{\Omega} |u(t)|^2 dx = \int_{\Omega} 2u'(t)u(t) dx$)
	\end{itemize}
\end{bem}
	
\mbox{} \\
Betrachte nun
\begin{equation} \left. \begin{cases}
	u_t + Lu = f & \text{ in } \Omega_T \\
	u = 0 & \text{ auf } \der \Omega \\
	u = y & \text{ auf } \setnull \times \Omega \end{cases} \right\} \text{ mit } f \in L^2(\Omega_T), g \in L^2(\Omega) \label{eq_A} \end{equation}
und $Lu = -\diver (A \nabla u) + b \cdot \nabla u + cu$, $A,b,c \in L^\infty(\Omega_T)$, $A$ symmetrisch positiv definit mit kleinstem Eigenwert $\lambda$. \\
Idee für schwache Lösung: \begin{itemize}
	\item Fasse $u$ auf als Funktion $u\colon [0,T] \rightarrow H_0^1(\Omega)$ (analog $f \colon [0,T] \rightarrow L^2(\Omega)$)
	\item Definiere Bilinearform $B_t[u,v] := \int_{\Omega} \nabla u \cdot A \nabla v + b \cdot \nabla uv + cuv dx = \int_{\Omega} Lu \cdot v dx$ für alle $u,v \in H_0^1(\Omega)$.
	\item Multipliziere $u_t + Lu = f$ mit $v \in H_0^1(\Omega)$ und integriere partiell \\
	$\Rightarrow (u_t(t),v)_{L^2(\Omega)} + B_t[u(t),v] = (f(t),v)_{L^2(\Omega)}$
	\item außerdem $u_t = g^0 + \sum_{i=1}^{n} g_{x_i}^i$ mit $g^0 = f - b \cdot \nabla u - cu, g^i = \sum_{i=1}^{n} A_{ij} \cdot u_{x_i}$, d.h.
	\[ ||u_t||_{H^{-1}(\Omega)} \leq \sqrt{ \sum_{i=0}^{n} ||g^i||_{L^2}^2} \leq C( ||u||_{H_0^1(\Omega)^t} ||f||_{L^1(\Omega)}) \]
\end{itemize}
$u \in L^2((O,T);H_0^1(\Omega))$ mit $u' \in L^2((0,T);H^{-1}(\Omega))$ heißt \Index{schwache Lösung} von \eqref{eq_A}, wenn $\langle u',v \rangle + B_t[u,v] = (f,v)_{L^2}$ für jedes $v \in H_0^1(\Omega)$ und für alle $t \in [0,T]$ sowie $u(0) = g$. (macht Sinn, da $u \in C([0,T];L^2(\Omega))$)

\begin{defn}[Garlerkin-Approximation]
Sei $\Omega$ Lipschitz. $L^2(\Omega)$ und $H_0^1(\Omega)$ sind Hilberträume. Man kann zeigen (Evans, Kapitel 6.5.1), dass glatte Funktionen $w_k \in H_0^1$ für $k \geq 1$ existieren, sodass gilt:
\begin{enumerate}[(i)]
	\item $(w_k)_k$ ist eine Orthogonalbasis von $H_0^1(\Omega)$.
	\item $(w_k)_k$ ist eine Orthonormalbasis von $L^2(\Omega)$.
\end{enumerate}
Für festes $m \in \NN$ approximieren wir nun eine schwache Lösung in $\langle w_1,w_2,\dots \rangle$, d.h. wir suchen
\[ u_m(t) = \sum\limits_{k=1}^{m} d_k^m(t) w_k \]
mit Koeffizienten $d_k^m \colon [0,T] \rightarrow \RR$, sodass
\begin{equation}
	\left. \begin{cases}
	(u_m',w_k)_{L^2} + B_t [u_m,w_k] = (f,w_k)_{L^2} & \forall t \in [0,T], k = 1,\dots,m \\ 
	\underbrace{(u_m(0),w_k)_{L^2}}_{d_m^k (0)} = (g,w_k) & \forall k
	\end{cases} \right\} \label{eq_B}
\end{equation}
\end{defn}

\begin{thm}[Existenz der Galerkin-Approximation]
	Für jedes $m \in \NN$ existiert eine eindeutige Lösung $u_m = \sum_{k=1}^{m} d_k^m w_k$ von \eqref{eq_B}
\end{thm}
	
\minisec{Beweis}
	Sei $B_t[w_l,w_k] = e_{kl}(t), D_m(t) = (d_m^1(t), \dots, d_m^m(t))^T, E(t) = (e_{kl}^{(i)})_{k,l = 1,\dots,m} \in \RR^{m \times m}$ sowie \linebreak
	$F(t) = ( (f(t),w_1)_{L^2},\dots,(f(t),w_m)_{L^2} )^T$. \\
	\eqref{eq_B} $\Leftrightarrow D_m'(t) + E(t) D_m(t) = F(t)$ bzw. $D_m'(t) = -E(t) D_m(t) + F(t)$ mit $D_m(0) = ( (g,w_1)_{L^2}, \dots, (g,w_m)_{L^2})^T$.\\
	Beachte: rechte Seite ist für jedes $t$ lipschitzstetig in $D_m$. Mit einem tiefen Satz für gewöhnliche Differentialgleichungen folgt, dass eine eindeutige Lösung $D_m(t)$ existiert. \qed
	
\mbox{} \\
Wie bei den elliptischen Gleichungen wollen wir annehmen, dass $\lambda, \Lambda, \gamma > 0$ existieren mit: \begin{itemize}
	\item $\xi' A(t,x) \xi \geq \lambda |\xi|^2$
	\item $| \xi^TA \xi| \leq \Lambda |\xi||\xi|$
	\item $\lambda^{-2} |b(t,x)|^2 + \lambda^{-1} |c(t,x)|^2 \leq \gamma^2, c(t,x) \geq 0$
\end{itemize}

\begin{thm}[Energieabschätzung]
\label{energieabsch}
	Es existiert ein $C < 0$ unabhängig von $m$, sodass:
	\[ \max_{t \in [0,T]} ||u_m(t)||_{L^2(\Omega)} + ||u_m||_{L^2((0,T);H_0^1(\Omega))} + ||u'_m||_{L^2((0,T);H^{-1}(\Omega))} \leq C \enbrace*{||f||_{L^2((0,T);L^2(\Omega))} + ||g||_{L^2(\Omega)}} \]
\end{thm}
	
\minisec{Beweis}
	Multipliziere \eqref{eq_B} mit $d_k^m$ und summiere $\sum_{k=1}^m \quad \Rightarrow \quad (u_m',u_m)_{L^2} + B_t[u_m,u_m] = (f,u_m)_{L^2}$. \\
	Wir hatten gezeigt (Theorem \ref{thm_71} bzw. Übungsblatt 9):
	\[ \beta ||u_m||_{H_0^1(\Omega)}^2 \leq B_t[u_m,u_m] + \gamma ||u_m||_{L^2(\Omega)}^2 \text{ für Konstanten } \beta,\gamma > 0 \]
	Damit folgt:
	\begin{equation}
	\begin{aligned}
	\underbrace{\frac{d}{dt} \frac{||u_m||_{L^2(\Omega)}^2}{2}}_{= \langle u_m',u_m \rangle} + \beta ||u_m||_{H_0^1(\Omega)}^2 &\leq \overbrace{\langle u_m',u_m \rangle + B_t[u_m,u_m]}^{=(f,u_m)_{L^2(\Omega)} = \int_{\Omega} \overbrace{f u_m}^{\leq \frac{1}{2} f^2 + \frac{1}{2} u_m^2} dx} + \gamma ||u_m||_{L^2(\Omega)}^2 \\
	&\leq \frac{1}{2} ||f||_{L^2(\Omega)}^2 + \enbrace*{\gamma + \frac{1}{2}} ||u_m||_{L^2(\Omega)}^2 \label{eq_C}
	\end{aligned}
	\end{equation}
	Mit $\eta(t) := ||u_m(t)||_{L^2(\Omega)}^2, \xi(t) := \frac{1}{2} ||f(t)||_{L^2(\Omega)}^2$ ergibt sich $\eta'(t) \leq \xi(t) + (\gamma + \frac{1}{2}) \eta(t)$. Mit Gronwall \ref{gronwall} folgt:
	\[ \eta(t) \leq e^{(\gamma + \frac{1}{2})t} (\eta(0) + \underbrace{\int_{0}^{t} \xi(s) ds}_{= \frac{1}{2} ||f||_{L^2((0,T);L^2(\Omega))}}), \]
	wobei $\eta(0) = (u_m(0),u_m(0))_{L^2} = \sum\limits_{k=1}^{m} |d_k^m(0)|^2 \leq \sum\limits_{k=1}^{\infty} |d_k^m(0)|^2 = ||g||_{L^2}^2$.
	\[ \Rightarrow \max_{t \in [0,T]} ||u_m||_{L^2} \leq C \sqrt{||g||_{L^2}^2 + ||f||_{L^2((0,T);L^2(\Omega))}} \leq \widetilde{C} (||g||_{L^2} + ||f||_{L^2((0,T);L^2(\Omega))}) \]
	Integriere nun \eqref{eq_C} von 0 bis $T$:
	\[ \frac{||u_m(T)||_{L^2}^2 - ||u_m(0)||_{L^2}^2}{2} + \beta ||u_m||_{L^2((0,T);H_0^1(\Omega))}^2 \leq \frac{||f||_{L^2((0,T);L^2(\Omega))}^2}{2} + (\gamma + \frac{1}{2}) \underbrace{\int_{0}^{T} ||u_m||_{L^2(\Omega)}^2 dt}_{\leq T \max_{t \in [0,T]} ||u_m||_{L^2(\Omega)}^2} \]
	\[ \Rightarrow ||u_m||_{L^2((0,T);H_0^1(\Omega))}^2 \leq C (||f||_{L^2((0,T);L^2(\Omega))}^2 + ||g||_{L^2(\Omega)}^2) \]
	Nun sei $v \in H_0^1(\Omega)$ beliebig mit $||v||_{H_0^1(\Omega)} \leq 1$.
	\begin{equation}
	\begin{aligned}
		\langle u_m',v \rangle &= (u_m',v)_{L^2} \overset{\text{ONB}}{=} \enbrace*{u_m', \sum\limits_{k=1}^\infty (v,w_k)w_k}_{L^2} = \enbrace*{u_m', \overbrace{\sum\limits_{k=1}^{m} (v,w_k)w_k}^{=: v_m}}_{L^2} \\
		&\overset{\eqref{eq_B}}{=} -B_t [u_m,v_m] + (f,v_m)_{l^2} \leq C (||u_m||_{H_0^1(\Omega)} + ||f||_{L^2(\Omega)}) \underbrace{||v_m||_{H_0^1(\Omega)}}_{\leq ||v||_{H_0^1(\Omega)} \leq 1} \notag
	\end{aligned}
	\end{equation}
	\[ \Rightarrow ||u_m'||_{H^{-1}(\Omega)} \leq C( ||u_m||_{H_0^1(\Omega)} + ||f||_{L^2}) \Rightarrow \int_{0}^{T} ||u_m'||_{H^{-1}(\Omega)}^2 dt \leq 2C (||u_m||_{L^2((0,T);H_0^1(\Omega))}^2 + ||f||_{L^2((0,T);L^2(\Omega))}) \]
	\qed
	
\begin{lemma}[Lemma von Gronwall] \label{gronwall}
	Sei $\eta(t) \geq 0$ absolut stetig \marginnote{4. Jul} (d.h. $\eta(t) = \int_{t_0}^{t} \dots ds$) mit $\eta'(t) \leq \Phi(t) \eta(t) + \Psi(t)$ fast überall für $\Phi,\Psi \geq 0$, $\int_{\RR} |\Phi|+|\Psi| dt < \infty$. Dann gilt:
	\[ \eta(t) \leq e^{\int_0^t \Phi(s) ds} \left\lbrack \eta(0) + \int_{0}^{t} \Psi(s) ds\right\rbrack \]
\end{lemma}

\minisec{Beweis}
	\[ \frac{d}{ds} (\eta(s) e^{-\int_{0}^{s} \Phi(r) dr} = e^{-\int_{0}^{s} \Phi(r) dr} \underbrace{(\eta'(s) - \eta(s) \Phi(s))}_{\leq \Psi(s)} \Rightarrow \eta(t) e^{-\int_{0}^{s} \Phi(r) dr} \leq \eta(0) + \int_{0}^{t} \underbrace{\Psi(s) e^{-\int_{0}^{s} \Phi(r)dr}}_{\leq \Psi(s)} ds \] \qed
	
\begin{thm}[Existenz schwacher Lösungen]
	\eqref{eq_A} besitzt eine schwache Lösung.
\end{thm}
	
\minisec{Beweis}
	\begin{enumerate}[1)]
		\item \begin{itemize}
			\item $u_m$ und $u_m'$ sind wegen \ref{energieabsch} gleichmäßig beschränkt in $X:= L^2((0,T);H_0^1(\Omega))$ bzw. $L^2((0,T);H^{-1}(\Omega))$.
			\item $X$ ist Hilbertraum mit Skalarprodukt $(u,v) = \int_{0}^{T} \int_{\Omega} uv + \nabla u \cdot \nabla v dx dt$ \\
			$\xLongrightarrow{\ref{riesz}}$ jedes Element $f \in X^*$ kann mit einem $u_f \in X$ identifiziert werden mit $||f||_{X^*} = ||u_f||_X$. \\
			$\Rightarrow X \equiv X^* \Rightarrow X^* \equiv X^{**} \Rightarrow X$ ist reflexiv. \marginnote{$\equiv$: isometrisch isomorph}
			\item $L^2((0,T),H^{-1}(\Omega)$ ist auch reflexiv (gleiches Argument, da $H^{-1}(\Omega) = (H_0^1(\Omega))^* \equiv H_0^1(\Omega)$.\\
			$\Rightarrow u_m \overset{L^2((0,T);H_0^1(\Omega)}{\rightharpoonup} u$ für Teilfolge und $u_m' \overset{L^2((0,T);H^{-1}(\Omega)}{\rightharpoonup} \widetilde{w}$ für Teilfolge. \\
			Weiterhin $\widetilde{w} = u'$, denn sei $\Phi \in C_c^\infty((0,T);\RR), w \in H_0^1(\Omega)$, dann ist
			\[ \underbrace{\langle \int_{0}^{T} u_m \Phi' dt, w \rangle}_{=\int_{0}^{T} \langle u_m, \Phi' w \rangle dt \xlongrightarrow{m \rightarrow \infty} \int_{0}^{T} \langle u,\Phi' w \rangle dt = \langle \int_{0}^{T} u \Phi' dt,w \rangle} = \underbrace{\langle - \int_{0}^{T} u_m' \Phi dt,w \rangle}_{= - \int_{0}^{T} \langle u_m',\Phi w\rangle dt \xlongrightarrow{m \rightarrow \infty} - \int_{0}^{T} \langle \widetilde{w},\Phi w \rangle dt = \langle -\int_{0}^{T} \widetilde{w} \Phi dt, w \rangle}  \]
			$\Rightarrow \int_{0}^{T} u \Phi' dt = -\int_{0}^{T} \widetilde{w} \Phi dt$
		\end{itemize}
		\item Sei $v(t) = \sum_{k=1}^{N} d_k(t) w_k$ für ein $N \in \NN$, da glatt, $N < m$. Aus \eqref{eq_B} folgt (multipliziere mit $d_k$, summiere, integriere):
		\begin{equation}
			\int_{0}^{T} \langle u_m',v \rangle dt + \int_{0}^{T} B_t [u_m,v] dt = \int_{0}^{T} (f,v)_{L^2} dt \label{eq_gronwall_stern}
		\end{equation}
		und mit $m \rightarrow \infty$ folgt weiter:
		\begin{equation}
			\int_{0}^{T} \langle u',v \rangle dt + \int_{0}^{T} B_t [u,v] dt = \int_{0}^{T} (f,v)_{L^2} dt \label{eq_gronwall_stern2}
		\end{equation}
		Dies gilt für alle $v \in L^2((0,T);H_0^1(\Omega))$, da sich diese durch $\sum_{k=1}^{N} d_k(t) w_k$ approximieren lassen, d.h.
		\[ \int_{0}^{T} d(t) (\langle u',v \rangle + B_t[u,v] - (f,v)_{L^2})dt = 0 \quad \text{für alle } v \in H_0^1(\Omega), d(t) \text{ messbar} \]
		\[ \Rightarrow \langle u',v \rangle + B_t[u,v] = (f,v)_{L^2} \text{ für fast alle } t \in [0,T] \]
		\item Sei $v \in C^1([0,T];H_0^1(\Omega))$ mit $v(T) = 0$. \eqref{eq_gronwall_stern} und \eqref{eq_gronwall_stern2} liefern:
		\[ \int_{0}^{T} -\langle u_m, v' \rangle dt + \int_{0}^{T} B_t [u_m,v] dt = \int_{0}^{T} (f,v)_{L^2} dt + (u_m(0),v(0)) \]
		\[ \int_{0}^{T} -\langle u, v' \rangle dt + \int_{0}^{T} B_t [u,v] dt = \int_{0}^{T} (f,v)_{L^2} dt + (u(0),v(0)) \]
		Für $m \rightarrow \infty$ folgt aus der ersten Gleichung:
		\[ \int_{0}^{T} -\langle u, v' \rangle dt + \int_{0}^{T} B_t [u,v] dt = \int_{0}^{T} (f,v)_{L^2} dt + (g,v(0)) \Rightarrow u(0) = g\] \qed
	\end{enumerate}

\begin{thm}[Eindeutigkeit schwacher Lösungen]
	Eine schwache Lösung von \eqref{eq_A} ist eindeutig.
\end{thm}
	
\minisec{Beweis}
	Seien $u_1, u_2$ schwache Lösungen, dann erfüllt $u := u_1 - u_2$:
	\[ \begin{cases}
		\langle u',v \rangle + B_t[u,v] = 0 & \text{ für alle } v \in H_0^1(\Omega) \text{ und fast alle } t \in [0,T] \\
		u(0)=0 \end{cases} \]
	Für $v = u$ ergibt sich
	\[ 0 = \frac{d}{dt} \enbrace*{\frac{||u||_{L^2(\Omega)}^2}{2}} + \underbrace{B_t[u,u]}_{\geq \beta ||u||_{H_0^1(\Omega)}^2 - \gamma ||u||_{L^2(\Omega)}^2} \]
	d.h. $\frac{d}{dt} ||u||_{L^2}^2 \leq 2\gamma ||u||_{L^2}^2$. Mit dem Lemma von Gronwall \ref{gronwall} folgt $||u||_{L^2}^2 \leq e^{2 \gamma t} ||u(0)||_{L^2}^2 = 0$. \qed
\newpage			% Schwache Lösungen von parabolischen Gleichungen
\section{Hyperbolische PDGL -- Wellengleichung}
\label{sec:para13}
	\[ u_{tt} - \Delta u = 0 \]
	physikalishe Interpretation: elastische Schwingung eines Materialstücks $\Omega$.
	\begin{itemize}
		\item Verschiebungsvektor $u$
		\item Dichte $\rho$
		\item Beschleunigung von $V \subset \Omega$: $\frac{d^2}{dt^2} \int_V udx = \int_V u_{tt} dx$
		\item auf $V$ wirkende elastische Kraft: $-\int_{\der V} F \cdot \nu dx = - \int_V \diver F dx$ ($F \in \RR^{n \times n}$: Spannungstensor)
		\item Spannungstensor $F \approx a Du$
		\item Newtons Gesetz der Bewegung: $\rho \frac{d^2}{dt^2} u dx = - \int_{\der V} F \cdot \nu dx \Rightarrow \rho u_{tt} = a \Delta u$
	\end{itemize}
	Physikalisch (und somit hoffentlich mathematisch) sinnvolle Randbedingungen:
	\[ \begin{cases}
	u_{tt} - \Delta u = 0	&	\text{auf } (0,T) \times \Omega	\\
	u = g	&	\text{auf } \setnull \times \Omega	\\
	u_t = h	&	\text{auf } \setnull \times \Omega \\
	u = f & \text{auf } \der \Omega
	\end{cases} \]
	
\begin{thm}[Explizite Lösung im Eindimensionalen]
	\begin{itemize}
		\item $\Omega = \RR$
		\item $0 = u_{tt} - \Delta u = u_{tt} - u_{xx} = (\der_t + \der_x)\overbrace{(\der_t - \der_x)u}^{=: v}$ \\
		$\Rightarrow v_t + v_x = 0 \Rightarrow v(t,x) = a(x-t) \Rightarrow u_t - u_x = a(x-t)$\\
		$\Rightarrow u(t,x) = \int_{0}^{t} a(x+(t-s)-s)ds + b(x+t) = \frac{1}{2} \int_{x-t}^{x+t} a(y) dy + b(x,t)$
		\item $u \big|_{t=0} = g = b, u_t \big|_{t=0} = a(x) + b'(x) = h \Rightarrow a = h-g'$ \\
		\[\Rightarrow u(t,x) = \frac{1}{2} \int_{x-t}^{x+t} h(y) - g'(y) dy + g(x+t) = \frac{g(x+t)+g(x-t)}{2} + \frac{1}{2} \int_{x-t}^{x+t} h(y) dy\]
	\end{itemize}
\end{thm}
	
\begin{bem}
	\begin{itemize}
		\item $u(t,x) = F(x+t) + G(x-t)$ für bestimmte $F,G$. Jede Funktion dieser Form löst $u_{tt}-u_{xx} = 0$.
		\item Für $g \in C^2(\RR), h \in C^1(\RR)$ ist $u \in C^2$, also tatsächlich eine klassische Lösung.
		\item Interpretation: $g(x+t)$: Welle, die nach links läuft, $g(x-t)$: Welle nach rechts.
	\end{itemize}
\end{bem}
	
\mbox{} \\
Auf der Halbachse:
	\[ \begin{cases}
	u_{tt} - \Delta u = 0	&	\text{auf } (0,T) \times (0,\infty)	\\
	u = g	&	\text{auf } \setnull \times (0,\infty)	\\
	u' = h	&	\text{auf } \setnull \times (0,\infty) \\
	u = 0 & \text{auf } \setnull \times \setnull
	\end{cases} \]
Reflektionsmethode: Definiere
\[ \widetilde{u}(t,x) = \begin{cases}
	u(t,x) & x \geq 0 \\
	-u(t,-x) & x < 0 \end{cases} \]
Analog für $g$ und $h$
\[ \Rightarrow \begin{cases}
	\widetilde{u}_{tt} - \Delta \widetilde{u} = 0 & \text{ auf } (0,\infty) \times \RR \\
	\widetilde{u} = g & \text{ auf } \setnull \times \RR \\
	\widetilde{u}_t = h & \text{ auf } \setnull \times \RR \end{cases} \]
\[ \Rightarrow \widetilde{u}(t,x) = \frac{1}{2} (\widetilde{g}(x+t) + \widetilde{g}(x-t)) + \frac{1}{2} \int_{x-t}^{x+t} h(y) dy \]
Interpretation: Wellen nach links werden bei 0 reflektiert.
\newpage			% Hyperbolische PDGL
\appendix							% Anhang beginnt hier (Nummerierung mit Buchstaben)
\section{Zusatznotizen}
\label{sec:SuppNotes}
	Sei $x \in \RR^n$ und $r > 0$. Bezeichne im Folgenden mit $B_r(x)$ die offene Kugel um $x$ mit Radius $r$, also
	\begin{equation}
		B_r(x) = \{ y \in \RR^n : |y-x| < r \}, \label{supp_0.1}
	\end{equation}
	und mit $\der B_r(x)$ den Rand von $B_r(x)$, also
	\begin{equation}
		\der B_r(x) = \{ y \in \RR^n : |y-x| = r \}. \label{supp_0.2}
	\end{equation}
	Analog bezeichnen wir mit $B_r$ die offene Kugel um den Ursprung und mit $\der B_r$ entsprechend ihren Rand. \\

\minisec{Die Coarea-Formel und Polarkoordinaten}
	Es gilt die folgende Integrationsregel:
	
\begin{thm} \label{supp_1.1}
	Sei $f \colon \RR^n \rightarrow \RR$ stetig und integrierbar. Dann gilt für alle $x_0 \in \RR^n$:
	\begin{equation}
		\int_{\RR^n} f(x) dx = \int_{0}^{\infty} \enbrace*{ \int_{\der B_r(x_0)}  f(x) dS(x) } dr, \label{supp_0.3}
	\end{equation}
	wobei $S$ das Oberflächenmaß auf $\der B_r(x_0)$ bezeichnet.
\end{thm}
	
\mbox{} \\
Theorem \ref{supp_1.1} kann mithilfe von Polarkoordinaten bewiesen werden und ist ein Spezialfall des folgenden Satzes:

\begin{thm}[Coarea-Formel] \label{supp_1.2}
	Sei $u \colon \RR^n \rightarrow \RR$ lipschitzstetig und für fast alle $r \in \RR$ sei die Niveaumenge $\{x \in \RR^n : u(x) = r\}$ eine glatte, $(n-1)$-dimensionale Hyperfläche des $\RR^n$. Ferner sei $f \colon \RR^n \rightarrow \RR$ stetig und integrierbar. Dann gilt: \index{Coarea-Formel}
	\begin{equation}
		\int_{\RR^n} f(x) |\nabla u(x)| dx = \int_{-\infty}^{\infty} \enbrace*{ \int_{ \{u = r\}} f(x) dS(x) } dr. \label{supp_0.4}
	\end{equation}
\end{thm}

\mbox{} \\
Die Coarea-Formel ist eine Art \glqq gekrümmte\grqq \ Version des Satzes von Fubini und ermöglicht die Umformung von $n$-dimensionalen Integralen in Integrale über Niveaumengen einer geeigneten Funktion.

\begin{bem} \label{supp_1.3}
	Theorem \ref{supp_1.1} folgt aus Theorem \ref{supp_1.2} für $u(x) = |x-x_0|$.
\end{bem}

\minisec{Kugelvolumen und Oberfläche}
	Um das Volumen $|B_r(x)|$ und das Maß $S(\der B_r(x))$ der Kugeloberfläche bestimmen zu können, benötigen wir zunächst die \Index{Gamma-Funktion}. Für $t > 0$ sei:
	\begin{equation}
		\Gamma(t) := \int_{0}^{+\infty} e^{-x} x^{t-1} dx. \label{supp_0.5}
	\end{equation}
	Wir prüfen zunächst, ob $\Gamma$ wohldefiniert ist: Sei $f(x) := e^{-x} x^{t-1}$. Dann ist $f(x) < x^{t-1}$ falls $x > 0$ und $t-1 > -1$. Da andererseits $\lim\limits_{x \rightarrow +\infty} x^{t+1}e^{-x} = 0$, existiert ein $M > 0$, sodass $x^{t+1}e^{-x} < 1$ für alle $x > M$, und damit $f(x) < \frac{1}{x^2}$ für alle $x > M$. Daraus folgt die Integrierbarkeit von $f$. Also ist $\Gamma(t) < \infty$ für alle $t > 0$. \\
	Die beiden Eigenschaften
	\begin{enumerate}[(1)]
		\item $\Gamma(1) = \int_{0}^{+\infty} e^{-x} dx = 1$
		\item $\Gamma(t+1) = \int_{0}^{+\infty} x^t e^{-x} dx = t \Gamma(t)$ für alle $t \geq 0$.
	\end{enumerate}
	zeigen, dass die Gamma-Funktion die Fakultätsfunktion auf $(0,\infty)$ fortsetzt; tatsächlich gilt für alle $n \in \NN$:
	\[ \Gamma(n+1) = n \Gamma(n) = n(n-1)\Gamma(n-1) = \dots = n!\Gamma(1) = n!\]
	Ein anderer Ausdruck für die Gamma-Funktion ist gegeben durch
	\begin{equation}
		\Gamma(t) = 2^{1-t} \int_0^{+\infty} e^{-\frac{y^2}{2}}y^{2t-1} dy, \label{supp_0.6}
	\end{equation}
	welchen man durch die Substitution $x = \frac{y^2}{2}$ erhält. \\
	Bezeichnen wir mit $Q_1 := [0,+\infty) \times [0,+\infty)$ den ersten Quadranten in der Ebene, folgt aus \eqref{supp_0.6} mithilfe von Fubini und Polarkoordinaten:
	\begin{equation}
	\begin{aligned}
		\left\lbrack \Gamma \enbrace*{\frac{1}{2}} \right\rbrack^2 &= 2 \enbrace*{ \int_0^{+\infty} e^{-\frac{x^2}{2}} dx} \enbrace*{ \int_{0}^{+\infty} e^{-\frac{y^2}{2} dy}} \\ \notag
		&= 2 \iint_{Q_1} e^{-\frac{x^2+y^2}{2}} dxdy = 2 \int_{0}^{\frac{\pi}{2}} d\vartheta \int_0^{\infty} \rho e^{-\frac{\rho^2}{2}} d\rho = \pi
	\end{aligned}
	\end{equation}
	und damit $\Gamma(1/2) = \sqrt{\pi}$. \\
	Sei nun $\omega_n$ das Volumen der Einheitskugel $B_1 \subset \RR^n$ und $\sigma_n$ das Maß der Oberfläche von $B_1$, d.h.
	\begin{equation}
		|B_r(x)| = r^n \omega_n, \qquad S(\der B_r(x)) = r^{n-1} \sigma_n \label{supp_0.7}
	\end{equation}
	
\begin{thm} \label{supp_2.1}
	Für $n \geq 2$ ist $\sigma_n = n \omega_n$.
\end{thm}
	
\minisec{Beweis}
	Mit $f \equiv 1$ folgt unmittelbar aus \eqref{supp_0.3}:
	\[	\omega_n = \int_{B_1} dx = \int_{0}^{1} \enbrace*{ \int_{\der B_\rho} dS} d\rho = \int_{0}^{1} S(\der B_\rho) d\rho = \sigma_n \int_0^1 \rho^{n-1} d\rho = \frac{\sigma_n}{n} \] \qed
	
Dies liefert und schließlich das Volumen $\omega_n$ der Einheitskugel für jedes $n \in \NN$:

\begin{thm} \label{supp_2.2}
	Sei $n \geq 1$, dann gilt:
	\begin{equation} 
		\omega_n = \frac{\pi^{n/2}}{(n/2) \Gamma (n/2)} \label{supp_0.8} 
	\end{equation}
\end{thm}
	
\minisec{Beweis}
	Mit $\Gamma(1) = 1$ und $\Gamma(1/2) = \sqrt{\pi}$ folgt, dass die Formel für $n=1$ und $n=2$ wahr ist:
	\[ \omega_1 = \frac{\pi^{1/2}}{(1/2)\Gamma(1/2)} = 2, \qquad \omega_2= \frac{\pi}{1 \cdot \Gamma(1)} = \pi \]
	Für $n \geq 3$ beweisen wir mit Induktion. Die Formel gelte also für $n - 2$, wobei $n \geq 3$. Dann gilt sie auch für $n$, denn: \\
	Für $x \in B_1 \subset \RR^n$ schreibe $x = (x',x'')$ mit $x' = (x_1,x_2)$ und $x'' = (x_3,\dots,x_n)$, dann ist
	\[ x' \in D_1 := \{(x_1,x_2) \in \RR^2 : x_1^2 + x_2^2 < 1 \} \]
	und
	\begin{equation}
	\begin{aligned}
		x'' \in (B_1)_{x'} := \ &\{x'' \in \RR^{n-2} : (x',x'') \in B_1 \} \\ \notag
			= \ &\{(x_3,\dots,x_n) \in \RR^{n-2} : x_3^2 + \dots + x_n^2 < 1 - x_1^2 - x_2^2 \}
	\end{aligned}
	\end{equation}
	Mit Fubini und der Induktionsannahme folgt:
	\begin{equation}
	\begin{aligned}
		\omega_n &= \int_{B_1} dx = \int_{D_1} dx' \int_{(B_1)_{x'}} dx'' 
			= \int_{D_1} \omega_{n-2} (1-x_1^2-x_2^2)^{(n-2)/2} dx_1 dx_2 \\
			&= \omega_{n-2} \int_0^{2\pi} d \vartheta \int_0^1 (1-\rho)^{(n-2)/2} \rho d\rho 
			= \frac{2\pi}{n} \omega_{n-2} = \frac{2\pi}{n} \frac{\pi^{(n-2)/2}}{\frac{n-2}{2} \Gamma \enbrace*{ \frac{n-2}{2}}} \\
			&= \frac{2\pi}{n} \frac{\pi^{(n-2)/2}}{\Gamma(n/2)} = \frac{\pi^{n/2}}{(n/2) \Gamma (n/2)} \notag
	\end{aligned}
	\end{equation}
	\qed
\newpage		% Zusatznotizen
%\newpage
\printindex			% Index ausgeben
\end{document}